% =====================================================================
%  OUTDATED -- describes the OLD implementation. Retained for reference only.
%  This draft documents the pre-DEFUN compiler: the typability-driven
%  specializer that selected one of three runtime targets (call-by-value /
%  call-by-name / interpret) and carved closed simply-typable sub-terms into
%  compiled "islands", plus the multi-stage island-depth self-compilation
%  climb. That whole architecture has been removed. The current compiler is
%  the DEFUN defunctionalization compiler (each Lam becomes an @interned
%  closure class, each App a Thunk); the self-compilation benchmark is now
%  the interpreter-vs-defunctionalized comparison in generated/defun-benchmark.
%  The live paper body is co-lambda.tex (input by preprint.tex / submission.tex);
%  this file is input by nobody. The "Shared body" note just below is itself
%  stale -- co-lambda.tex took that role. Do not treat the prose below as a
%  description of the present implementation.
% =====================================================================
%  Solving Coalgebraic Equations in the Pure $\lambda$-Calculus
% ---------------------------------------------------------------------
%  Shared body. This file carries NO \documentclass: it is \input by two
%  entry points that supply the class (mimicking ../inheritance-calculus):
%    * preprint.tex   -- arXiv / preprint  (acmart, nonacm)
%    * submission.tex -- POPL submission   (acmart, review, anonymous)
%  Build both with: latexmk  (see latexmkrc)
% ---------------------------------------------------------------------
%  STATUS: working draft.
%  * Placeholders are marked \TODO{...} and render in red.
%  * Every bibliography entry is UNVERIFIED: author/title are best-effort
%    recollections; years, venues, volumes, and pages are marked [verify].
%    Do not cite from this file without checking each entry.
%  * Proofs are sketches; load-bearing gaps are flagged with \TODO.
% =====================================================================

%% acmart already loads amsmath, amssymb, amsthm, xcolor, and hyperref;
%% only mathtools needs to be added here.
\usepackage{mathtools}

%% listings renders the machine-generated Python in the compiler appendix.
\usepackage{listings}
\lstset{
  basicstyle=\ttfamily\small,
  breaklines=true,
  breakatwhitespace=false,
  columns=fullflexible,
  keepspaces=true,
  showstringspaces=false,
  language=Python,
  xleftmargin=1em,
  aboveskip=0.4\baselineskip,
  belowskip=0.4\baselineskip,
}

% --- draft placeholder macro ------------------------------------------
\newcommand{\TODO}[1]{\textcolor{red}{[\textbf{TODO:} #1]}}
\newcommand{\citeTODO}[1]{\textcolor{red}{[\textbf{cite:} #1]}}

% --- general notation --------------------------------------------------
\newcommand{\Set}{\mathbf{Set}}
\newcommand{\Beh}{B}                         % behaviour functor B = M F
\newcommand{\rfp}{\varrho}                   % rational fixpoint
\newcommand{\behav}{\mathsf{beh}}            % behaviour map into the final coalgebra
\newcommand{\outm}{\mathsf{out}}             % structure map: the iteration of the equation morphism
\newcommand{\stepm}{e}                       % equation morphism (the map e)
\newcommand{\Delaym}{\mathsf{Delay}}
\newcommand{\Pow}{\mathcal{P}}
\newcommand{\iter}[1]{#1^{\dagger}}          % Elgot iteration

% --- lambda-instance notation (used only in the HNF bridge/examples) ---
\newcommand{\Sh}{\mathsf{Sh}}
\newcommand{\Lamsh}{\mathsf{Lam}}
\newcommand{\Appsh}{\mathsf{App}}
\newcommand{\Varsh}{\mathsf{Var}}
\newcommand{\dlam}{\mathord{\downarrow}\lambda}
\newcommand{\dhead}{\mathord{\downarrow}\mathsf{h}}
\newcommand{\darg}{\mathord{\downarrow}\mathsf{a}}
\newcommand{\subt}[2]{#1\!\downarrow\!#2}
\newcommand{\reduces}{\twoheadrightarrow}
\newcommand{\bred}{\rightarrow_{\beta}}
\newcommand{\present}{\mathtt{present}}

% --- theorem environments ---------------------------------------------
%% acmart provides theorem/lemma/proposition/corollary (acmplain) and
%% definition/example (acmdefinition), all sharing the `theorem' counter.
%% Only `remark' is missing, so we add it on the same shared counter.
\AtBeginDocument{%
  \theoremstyle{acmdefinition}%
  \newtheorem{remark}[theorem]{Remark}%
}

\begin{document}

\title{Solving Coalgebraic Equations in the Pure $\lambda$-Calculus}

\author{Bo Yang}
\affiliation{%
  \institution{Figure AI Inc.}
  \city{San Jose}
  \state{California}
  \country{USA}
}
\email{yang-bo@yang-bo.com}
\thanks{The author completed this work independently and outside the scope of employment at Figure AI. No Figure AI time, funding, or other resources were used in support of this work.}

\begin{abstract}
  Representing and transforming cyclic and infinite data in the pure $\lambda$-calculus normally requires an added recursion construct, a \texttt{letrec}, a $\mu$-binder, or a built-in $Y$ for graph reduction: such data has no finite normal form, so ordinary reduction unfolds it forever instead of folding it into a finite cyclic graph, and an ordinary recursive function over it does the same. We show no extension is needed, and the denotational semantics is unchanged: a term keeps its standard meaning, which we represent as a finite cyclic graph when it is rational. The insight is that such data is the unique solution of a coalgebraic equation, and that any coalgebraic equation can be solved directly, by tabling: when a state recurs, the solver folds it into a back edge, so rational data, that with finitely many distinct subtrees, comes out as a finite cyclic graph. The solver halts when finitely many states are reachable, within that rational fragment, and the rational fragment is optimal, since data beyond it has no finite representation, so no denotation-preserving procedure halts on more. The algorithm's first application is the pure $\lambda$-calculus itself: head normalisation is such an equation, and because its states are first-order graph nodes rather than functions their equality is decidable where $\beta$-equality is not, so the solver is an interpreter. It decides $\Omega$ as $\bot$ in finite time and evaluates $\mathtt{map}\,\mathtt{succ}\,(Y\,(\mathtt{cons}\;0))$, an ordinary map over the cyclic stream of zeros, to the cyclic stream of ones as a finite cyclic graph, with no construct added. Because any coalgebra Scott-encodes as a $\lambda$-term, this one interpreter solves them all, making the pure $\lambda$-calculus a domain-specific language whose memoization and transposition tables fall out of node identity: Datalog least Herbrand models, graph reachability and Andersen points-to analysis, minimax search, and dynamic programming. A second example coalgebra is compilation. Its equation sends each sub-term of a source program, by specialization rules, to a node of a Python AST over that sub-term's own sub-terms, and its solution is the program optimised to that AST. A compiler computes it by the same tabling, compiling each shared sub-term once into a finite graph, and is itself a $\lambda$-term the interpreter solves, so program analysis is expressed as evaluation, and it compiles its own source in a self-hosting bootstrap.
\end{abstract}

\keywords{coalgebra, coalgebraic equations, rational fixpoint, algebraic compactness,
  structure as identity, least fixpoint, rational trees, tabling, head reduction, Datalog,
  program analysis, domain-specific language,
  \texorpdfstring{$\lambda$-calculus}{lambda-calculus},
\texorpdfstring{B\"ohm trees}{Bohm trees}}

\maketitle

% =====================================================================
\section{Introduction}\label{sec:intro}
% =====================================================================
A coalgebra is a state set $X$ with a structure map that, at each state, exposes one
layer of behaviour and names the successor states. Its \emph{behaviour}\footnote{Throughout,
  \emph{behaviour} is meant in the coalgebraic sense: the tree a state unfolds to, that is, its
  denotation, and not the operational behaviour in the sense of a reduction sequence or contextual
  equivalence. In the $\lambda$-instance (Section~\ref{sec:bridge}) it is the L\'evy--Longo or B\"ohm
  tree; the reduction is the internal process the structure map collapses, while the behaviour is the
tree that process exposes. Two terms may thus share a behaviour yet reduce differently. This tree is the
  first-order datum over the signature $F$, not the extensional behaviour of a higher-order value; and
  behavioural \emph{equivalence}, sharing that tree, is a comparison among states, coarser than the
  structural node identity the tabling uses to fold cycles, and in the $\lambda$-instance is B\"ohm-tree
  equality, hence undecidable; we solve a single equation rather than compare states, so it is not a
  relation the construction computes (Section~\ref{sec:intro}).} is the
unfolding of that map: the possibly infinite tree obtained by exposing layer after layer. We make one
observation about this unfolding and pursue its consequences. When the behaviour is
\emph{rational},\footnote{Throughout, \emph{rational} qualifies a behaviour, that is, a tree: the
  tree has finitely many distinct subtrees, equivalently it is the unfolding of a finite graph (a
  regular, or Courcelle, tree). This is the rational-fixpoint sense, unrelated to rational numbers,
and it is a property of the behaviour, not of the term or its reduction.} with finitely many
distinct subtrees, it is the unfolding of a finite
cyclic graph, and that graph is computed by \emph{tabling}: identify a state with its own
behaviour, table the states seen so far, and when a state recurs while its behaviour is computed,
tie a back edge instead of unfolding forever.

Two things make this effective. The first is the form the input takes; the second is what makes its
cycles detectable at all.

First, the structure map need not be handed to us as a one-step observation. It may itself be
produced by an \emph{internal process}: from a state, perform some internal steps, possibly
none, until a layer of behaviour is exposed. We capture the internal process by an effect
monad $M$ and the exposed layer by a signature functor $F$, so the datum is a coalgebra
$\stepm\colon X \to M(F X + X)$ valued in the \emph{resumption} of $F$ over $M$
(Definition~\ref{def:itcoalg}): the $F X$ summand exposes a layer, the $X$ summand continues
silently, and $M$ carries the effect of the step. The one-step structure map
$\outm \colon X \to M(F X)$ is then recovered as the least solution of the iteration that runs
the silent steps (Definition~\ref{def:out}). Because it is a \emph{least} fixpoint, it ascends
from $\bot$, so a silent cycle, a state whose internal process returns to it with no layer
exposed, terminates as $\bot$ in finite time where the process alone diverges. Head
normalisation of the $\lambda$-calculus is exactly this situation, with weak-head reduction as
the internal process (Section~\ref{sec:bridge}).

Second, recognising that a state has recurred is a test of \emph{identity}. This is the word the
title turns on. The identity is decidable because the states are \emph{first-order}: they range
over graph nodes, never over a function space. The cycle fold is then nothing special; it is
ordinary tabling on state identity when the tabling core is the unfolding map. Where
the states are higher-order values, as in the na\"ive reading of the $\lambda$-calculus, the only
available identity is $\beta$-equality, which is undecidable, and neither memoisation nor cycle
folding is computable. Reading behaviour off a first-order coalgebra is what makes identity
decidable and the folding of cyclic structure effective.

\paragraph{Why this answers the cyclic-data question}
It is folklore that the pure $\lambda$-calculus cannot represent or manipulate cyclic data without
an added recursion construct (a \texttt{letrec}, a $\mu$-binder, observable sharing, mutable
references). The construction refutes this: the $\lambda$-calculus, read as an equation morphism,
is one instance among many, and on it a pure term such as $Y\,(\mathtt{cons}\;0)$ has a rational
behaviour, so tabling folds it into a finite circular graph and an ordinary recursive function over
it terminates, with nothing added to the calculus (Sections~\ref{sec:bridge}, \ref{sec:application}).
The capability is not special to the $\lambda$-calculus; it is a property of any equation morphism
with rational behaviour.

\paragraph{Contributions}
\begin{itemize}
  \item A reading of the behaviour of an equation morphism $\stepm \colon X \to M(F X + X)$ by
    tabling its one-step structure map over the reachable states, parametric in an effect monad
    $M$ (the internal process) and a signature functor $F$ (the exposed layer)
    (Sections~\ref{sec:domain}, \ref{sec:itcoalg}). The one-step map $\outm = \iter{\stepm}$ is
    the least solution of the silent iteration (Definition~\ref{def:out}), so divergence is
    decided as $\bot$ in finite time on the rational fragment.
  \item One fixpoint, and a graph for free (Sections~\ref{sec:itcoalg}, \ref{sec:readout}). The
    only fixpoint is the silent collapse $\outm = \iter{\stepm}$, an \emph{unguarded least}
    fixpoint resolving a silent cycle to $\bot$ (the bottom of $M$). Tabling $\outm$ over the
    reachable states yields the finite cyclic graph directly: its nodes are the states, and a
    cycle is automatic when $\outm$ names a state already tabled. No separate behaviour fixpoint is
    needed; the graph is the reachable sub-coalgebra $(R(M),\outm)$.
  \item Adequacy (Theorem~\ref{thm:adequacy}): the reachable sub-coalgebra $(R(M),\outm)$ unfolds to
    the behaviour of $\stepm$ up to silent steps, agreeing layer for layer with what the coalgebra
    enumerates. Termination (Theorem~\ref{thm:termination}): tabling halts iff finitely many
    states are reachable, and then the graph is a finite element of the rational fixpoint
    $\rfp\Beh$. Two independent finiteness conditions appear (Remark~\ref{rem:tworationalities}).
  \item Least equals greatest (Theorem~\ref{thm:unique}): because $\outm$ exposes a constructor or
    $\bot$, the behaviour functional is guarded, hence a contraction with a \emph{unique} fixpoint,
    so the least-fixpoint reading and the coinductive unfolding coincide and no least-versus-greatest
    choice arises. This is the operational face of algebraic compactness over the $\bot$-completed
    domain (Remark~\ref{rem:compact}); the least fixpoint is indispensable only at the silent
    collapse, where it decides divergence as $\bot$ (Remark~\ref{rem:leastessential}).
  \item A bridge (Section~\ref{sec:bridge}): head and weak-head normalisation of the pure
    $\lambda$-calculus are the instance $M = \Delaym$ (pure divergence) and $F$ the
    $\lambda$-shape signature. There $\outm$ is the weak-head shape relation, $\Omega$ reads out
    as $\bot$, and a guarded cyclic term reads out as a finite cyclic graph. Deterministic
    reduction strategies in general are the case $M = \Delaym$ (Remark~\ref{rem:determinism}).
  \item As an application, structural computation over cyclic data inside the pure
    $\lambda$-calculus (Section~\ref{sec:application}): an ordinary $\mathtt{map}$, with nothing
    cycle-aware in it, folds a circular list into a circular list, and a na\"ive walk terminates
    where head reduction diverges.
\end{itemize}

\paragraph{Supplementary material}
An executable implementation accompanies the paper: a first-order interpreter for the
$\lambda$-instance of Section~\ref{sec:bridge} ($\Varsh$/$\Lamsh$/$\Appsh$, no recursion binder), in
which a term's behaviour is read by tabling the weak-head structure map (Definition~\ref{def:shbridge}),
whose least-fixpoint silent collapse decides an unproductive cycle as $\bot$ and whose reachable
sub-coalgebra is the finite cyclic graph, printed with explicit back-references.\anon[ It is included
as supplementary material.]{ It is openly available
at \url{https://github.com/Atry/MIXINv2}~\cite{colambda}.} It realises the construction
directly: the silent collapse $\outm$ as a least fixpoint (with a finite-budget first iterate as a
policy flag), and the tabling that folds on structural identity, realised by interning so identity is
a pointer test. It also provides the head structure map of Remark~\ref{rem:boehm} as a second
least-fixpoint instance over the same carrier and behaviour map. The examples of
Sections~\ref{sec:overview}, \ref{sec:bridge}, and \ref{sec:application} are reproduced as tests:
$Y\,(\mathtt{cons}\;0)$ folds to a finite cyclic graph, $\Omega$ and $Y\,(\lambda x.x)$ read out as
$\bot$, the na\"ive walk terminates over the cyclic stream where head reduction loops, the
rationality dichotomy of Theorem~\ref{thm:termination} is exercised on a rational and a non-rational
stream, and the weak-head and head readings are checked to terminate together while differing on
$\lambda x.\,\Omega$ (Remark~\ref{rem:boehm}).

% =====================================================================
\section{Overview}\label{sec:overview}
% =====================================================================
Take the cyclic stream of zeros, $r = \mathtt{cons}\;0\;r$. As the behaviour of a coalgebra it is the
infinite tree $\mathtt{cons}\;0\;(\mathtt{cons}\;0\;\cdots)$, but it has only \emph{two} distinct
subtrees, $0$ and $r$ itself, so it is rational. Reading it off by tabling, the cell
$\mathtt{cons}\;0\;(-)$ is registered before its tail is followed; the tail returns to that very
cell, which is already in the table, so the tabling ties a back edge and stops. The result is a
finite graph, a node pointing to itself, that unfolds to the infinite stream. Any presentation of $r$
as an equation morphism with this behaviour gives the same finite graph, the $\lambda$-term
$Y\,(\mathtt{cons}\;0)$ among them (Section~\ref{sec:bridge}).

Now take a process that returns to a state \emph{without} ever exposing a layer: an
\emph{unproductive} cycle. There is nothing to tie a back edge to. This is the one place a fixpoint is
needed: the structure map $\outm$ is the least fixpoint of the silent iteration, so an unproductive
cycle, having nothing to fill its hole, settles at $\bot$. In the $\lambda$-instance this is reached
by a recursive binder ($\mathtt{letrec}\;x=x$) and equally by self-application,
$\Omega=(\lambda x.x\,x)(\lambda x.x\,x)$, which steps to itself with no constructor exposed, so
$\Omega = \bot$, decided in finite time where reduction loops forever.

So there is a single fixpoint, the silent collapse $\outm$. The guarded cycle above is not a second
fixpoint but a back edge the tabling ties for free: the tail already \emph{is} the registered cell, so
the cyclic graph is there in the coalgebra and computing the behaviour merely prints it. The guarded case folds,
the unproductive case bottoms out, and whether either terminates is governed by one property,
rationality, developed next.

% =====================================================================
\section{The behaviour domain}\label{sec:domain}
% =====================================================================
We fix the codomain of the behaviour map, the trees a coalgebra unfolds to, and the order along which the
fixpoints below ascend, before introducing the coalgebra that drives them.

\subsection{Signature and effect}
\begin{definition}[Signature functor]\label{def:signature}
  A \emph{signature} is a finitary container functor $F\colon \Set \to \Set$,
  $F X = \coprod_{\sigma \in \Sigma} X^{\mathit{ar}(\sigma)}$, for a set $\Sigma$ of operations with
  finite arities $\mathit{ar}(\sigma) \in \mathbb{N}$. We require equality of \emph{shapes} (operations
  together with their arity, i.e.\ the elements of $\Sigma$) to be decidable and hashable. The flat
  domain of shapes is $F_{\bot} X = F X \cup \{\bot\}$, ordered by $\bot \sqsubseteq c$ for every $c$
  and no other strict relation.
\end{definition}

\begin{definition}[Iteration-ready (Elgot) monad]\label{def:elgot}
  An \emph{iteration-ready} monad is a monad $M$ on $\Set$ such that each $M A$ is a pointed
  $\omega$-complete partial order with least element $\bot_{MA}$ (the \emph{divergent} computation), the
  unit $\eta$ and Kleisli extension $(-)^{*}$ are $\omega$-continuous, and $\bot$ is preserved by
  Kleisli composition (a divergent prefix diverges). Consequently every $f\colon X \to M(Y + X)$ has a
  least \emph{iteration} $\iter{f}\colon X \to M Y$, the least fixpoint of
  $g \mapsto \big(\,[\,\eta\circ\mathrm{inl},\,g\,]^{*}\big)\circ f$ in the pointwise order. This is the
  order-theoretic presentation of a \emph{complete Elgot monad} \citeTODO{adamek-milius-velebil,
  goncharov-elgot}.
\end{definition}

\begin{example}[Effects]\label{ex:effects}
  $M = \Delaym$, the delay monad $\Delaym A = \nu Z.\,(A + Z)$ \citeTODO{capretta}, is the free
  iteration-ready monad on the identity: $\bot$ is the infinite \texttt{later}, divergence. It models a
  \emph{deterministic} process with no effect but possible non-termination, and is the case the
  $\lambda$-calculus instance uses (Section~\ref{sec:bridge}). $M = \Pow$ (powerset) models
  non-determinism, with $\bot = \varnothing$ and the iteration a set-valued least fixpoint; $M$ a
  subdistribution monad models probabilistic choice; $M = S \Rightarrow (S \times -)$ models state.
  Every definition and theorem below is stated for an arbitrary iteration-ready $M$.
\end{example}

\begin{definition}[Behaviour functor]\label{def:behfun}
  The \emph{behaviour functor} is $\Beh = M\circ F$. A one-step coalgebra is a map $X \to \Beh X$. We
  write $\nu\Beh$ for the final $\Beh$-coalgebra (the domain of behaviours) and $\rfp\Beh$ for its
  \emph{rational fixpoint}, the colimit of the finite $\Beh$-coalgebras, equivalently the behaviours
  with finitely many distinct subtrees \citeTODO{adamek-milius-velebil-rational}.
\end{definition}

\subsection{Trees, order, rationality}
For the order and the explicit tree presentation we spell out the deterministic skeleton; with a
non-trivial $M$ each node additionally carries the $M$-structure of its layer (a set of layers for
$M=\Pow$, a distribution for a probabilistic $M$, and so on), and the order and completeness lift
$M$-pointwise.

\begin{definition}[Trees]\label{def:domain}
  A \emph{tree} is a partial $\Sigma$-labelling of the arity-respecting positions reachable by
  descent, with a distinguished $\bot$ leaf for an exposed-nothing node. Formally a tree is a map
  $t\colon \mathrm{dom}(t) \to \Sigma \cup \{\bot\}$ with $\mathrm{dom}(t)$ prefix-closed under the
  descent into children and arity-respecting: a node labelled $\sigma$ has exactly $\mathit{ar}(\sigma)$
  children and a node labelled $\bot$ has none. Let $D$ be the set of all finite or infinite such trees,
  carrying the $M$-structure of Definition~\ref{def:behfun} at each node.
\end{definition}

\begin{definition}[Approximation order]\label{def:order}
  For $t,u \in D$, $t \sqsubseteq u$ iff $\mathrm{dom}(t) \subseteq \mathrm{dom}(u)$ and for every
  $p \in \mathrm{dom}(t)$ either $t(p) = u(p)$ or $t(p) = \bot$. Thus $u$ refines $t$ by replacing
  $\bot$-leaves with subtrees (and, under a non-trivial $M$, by refining the $M$-structure
  pointwise).
\end{definition}

\begin{proposition}[The codomain is a domain]\label{prop:cpo}
  $(D,\sqsubseteq)$ is a pointed $\omega$-complete partial order with least element
  $\bot_D$ (the single $\bot$ leaf), and least upper bounds of $\sqsubseteq$-directed families exist and
  are computed nodewise. Proof in Appendix~\ref{app:domain}.
\end{proposition}

\begin{definition}[Rational behaviour]\label{def:rational}
  $t \in D$ is \emph{rational} if its set of subtrees is finite; equivalently $t$ is the unfolding of a
  finite $\Sigma$-labelled graph, equivalently $t \in \rfp\Beh$ (Definition~\ref{def:behfun}).
\end{definition}

% =====================================================================
\section{Equation morphisms and the structure map}\label{sec:itcoalg}
% =====================================================================
The development is parametric in the internal process that produces the structure map. We give it,
extract the one-step map by collapsing the silent steps, and record that the collapse is a least
fixpoint.

\begin{definition}[Equation morphism]\label{def:itcoalg}
  An \emph{equation morphism} is a map
  \[
    \stepm \colon X \longrightarrow M\big(F X + X\big),
  \]
  where $X$ is a set of \emph{states} with decidable, hashable equality. From a state, $\stepm$
  performs one $M$-step whose result is either a layer of behaviour, $\mathrm{inl}(t)$ with $t \in F X$,
  or a silent continuation, $\mathrm{inr}(x')$ with $x' \in X$. The $\mathrm{inl}$ alternative is
  \emph{guarded} (a layer is exposed), the $\mathrm{inr}$ alternative \emph{unguarded} (the process
  continues with no layer). Reachability is closed under both alternatives, so the children of an
  exposed layer and every silent successor are again states.
\end{definition}

\begin{definition}[The structure map]\label{def:out}
  The \emph{structure map} is the iteration that runs the silent steps to the next layer,
  \[
    \outm \;:=\; \iter{\stepm} \colon X \longrightarrow M(F X),
  \]
  the least fixpoint (Definition~\ref{def:elgot}) of
  $\Phi(g)(x) = \big(\,[\,\eta\circ\mathrm{inl},\,g\,]^{*}\big)(\stepm(x))$. Equivalently $\outm(x)$ runs
  $\stepm$ from $x$, following $\mathrm{inr}$ continuations, and returns the first $\mathrm{inl}$ layer
  under the accumulated $M$-effect; where the silent process never exposes a layer, $\outm(x)$ is the
  divergent computation $\bot_{M(FX)}$.
\end{definition}

\begin{proposition}[The structure map is well-defined]\label{prop:outwd}
  For every iteration-ready $M$ and every equation morphism $\stepm$, $\Phi$ is monotone and
  $\omega$-continuous in the pointwise order on $X \to M(F X)$, so by Kleene
  $\outm = \bigsqcup_{k} \Phi^{k}(\lambda x.\,\bot)$ exists and is the least fixpoint, and $\outm(x)$ is
  the divergent computation exactly when the silent process from $x$ exposes no layer. Proof in
  Appendix~\ref{app:out}.
\end{proposition}

\begin{remark}[Why the bottom is $\bot$, and why that is sound]\label{rem:bottom}
  The iteration ascends from a bottom value, and that bottom is the divergent computation $\bot$ of $M$,
  fresh and distinct from every exposed layer. This is what makes the resolution test exact: a state is
  resolved precisely when its value is no longer $\bot$, so a silent cycle is recognised as soon as it
  returns to a state already on the stack, and decided as $\bot$ in finite time. Reusing an exposed
  value as the bottom would make a still-unresolved state indistinguishable from one that genuinely
  computed that value, and a round could read a placeholder as a real value and reach a wrong fixpoint.
  The only identity test, the one that lets the iteration recognise a cycle, is the decidable equality
  of states, not an equality of the values they compute.
\end{remark}

The structure map turns the equation morphism $\stepm$ into an ordinary one-step coalgebra
$\outm \colon X \to \Beh X$ for $\Beh = M F$ (Definition~\ref{def:behfun}). The rest of the
development runs on $\outm$ and never re-inspects the silent steps. The next definition is the
positive presentation that makes $\Phi$ monotone above; like the rest, it is parametric in $M$ and
$F$.

\begin{remark}[Positivity]\label{rem:positivity}
  The clauses of $\stepm$ are \emph{positive}: the choice between exposing a layer and continuing is a
  positive disjunction over the $M$-result, never a negation-as-failure test such as ``no layer is
  exposable''. Positivity is what makes $\Phi$ monotone, hence $\outm$ a least fixpoint by
  Knaster--Tarski, so well-definedness of $\outm$ costs nothing beyond fixing the positive presentation.
  The work of the paper is elsewhere (Sections~\ref{sec:readout}--\ref{sec:application}).
\end{remark}

% =====================================================================
\section{The rational graph: the reachable sub-coalgebra}\label{sec:readout}
% =====================================================================
The structure map is already an ordinary one-step coalgebra $\outm \colon X \to \Beh X$
(Definition~\ref{def:out}) whose successors are again states, so nothing further need be built: the
graph is the part of $(X,\outm)$ reachable from the root. Computing the behaviour is tabling $\outm$.

\begin{definition}[Reachable sub-coalgebra]\label{def:reach}
  For a state $M$, the \emph{reachable set} $R(M) \subseteq X$ is the least set containing $M$ and
  closed under the successors of $\outm$: if $x \in R(M)$ and $\outm(x)$ exposes a layer
  $\sigma(x_1,\dots,x_k)$, then $x_1,\dots,x_k \in R(M)$ (for a branching $M$, every state in every
  alternative of the layer). The \emph{graph} of $M$ is the sub-coalgebra $(R(M),\outm)$: nodes are the
  reachable states, edges are $\outm$, and a node with $\outm(x)=\bot$ is the meaningless leaf.
\end{definition}

\begin{definition}[Reading the graph; the behaviour it denotes]\label{def:behreadout}
  The graph $(R(M),\outm)$ is computed by \emph{tabling} $\outm$: keep a table of resolved states; from
  $M$, compute $\outm(M)$ and, for each successor state, reuse its table entry or recurse and add one. A
  successor already tabled closes a \emph{cycle}; a state with $\outm(x)=\bot$ is a leaf. The
  \emph{behaviour} $\behav(M) \in \nu\Beh$ is what this graph denotes: the unfolding of $(R(M),\outm)$,
  equivalently the unique $\Beh$-coalgebra morphism $X \to \nu\Beh$ at $M$ (finality). The behaviour is
  never materialised; the graph is its finite presentation, and every question about it is answered on
  the graph.
\end{definition}

\begin{remark}[One fixpoint; cycles for free]\label{rem:twofix}
  There is only one fixpoint: the silent collapse $\outm = \iter{\stepm}$ (Definition~\ref{def:out}), an
  unguarded least fixpoint that decides an unproductive cycle as $\bot$. Building the graph is not a
  second fixpoint. A cycle appears the moment $\outm$ names a state already tabled, so the back edge is
  structural, not tied by a separate computation, and ``reading out'' is just enumerating
  $(R(M),\outm)$; printing it with explicit back-references is a traversal of that finite graph, again no
  fixpoint. The whole algorithm is: table $\outm$.
\end{remark}

\begin{definition}[Guarded and unproductive states]\label{def:guarded}
  A state $x$ is \emph{cyclic} if a state equal to it recurs strictly below it on some descent of
  $\outm$. A cyclic $x$ is \emph{guarded} if $\outm(x)$ exposes a layer before the re-entry, and
  \emph{unproductive} if $\outm(x) = \bot$, so the silent process of $x$ returns to it with no layer
  exposed. A guarded cycle is a graph node pointing back to an earlier node, the finite representation of
  a rational behaviour; an unproductive cycle is the $\bot$ leaf.
\end{definition}

% =====================================================================
\section{Adequacy and termination}\label{sec:adequacy}
% =====================================================================
\begin{theorem}[Adequacy]\label{thm:adequacy}
  For every equation morphism $\stepm \colon X \to M(F X + X)$, the graph $(R(M),\outm)$
  (Definition~\ref{def:reach}) unfolds to the behaviour of $\stepm$ up to silent steps: its unfolding
  $\behav(M)$ is the unique $\Beh$-coalgebra morphism $X \to \nu\Beh$ through the structure map
  $\outm = \iter{\stepm}$ (Definition~\ref{def:behreadout}), so the graph agrees layer for layer with
  the unfolding the coalgebra enumerates once the silent steps are run. Proof in
  Appendix~\ref{app:adequacy}.
\end{theorem}

\noindent
There is no normal-form translation to justify and no administrative bridge: once the silent steps are
collapsed into $\outm$, the graph is the behaviour's finite presentation essentially by construction,
and the only content is that the silent collapse $\iter{\stepm}$ presents the internal process
faithfully (Proposition~\ref{prop:outwd}).

\begin{theorem}[Termination]\label{thm:termination}
  For every equation morphism with decidable state equality, tabling $\outm$ from $M$ halts if and
  only if finitely many states are reachable from $M$ (the reachable set $R(M)$ is finite,
  Definition~\ref{def:reach}), and in that case the graph $(R(M),\outm)$ is a finite cyclic graph whose
  unfolding lies in the rational fixpoint $\rfp\Beh$ (Definition~\ref{def:rational}). When infinitely
  many states are reachable, the tabling enumerates them and diverges. Proof in
  Appendix~\ref{app:termination}.
\end{theorem}

\begin{remark}[Two finiteness conditions]\label{rem:tworationalities}
  Two independent conditions govern termination, both with respect to structural state identity. The
  silent collapse $\outm(x)$ (Definition~\ref{def:out}) halts iff the silent orbit from $x$ is finite,
  returning to a state already seen rather than producing ever-new ones. The graph
  (Definition~\ref{def:behreadout}) halts iff the reachable states $R(M)$ are finite. They are logically
  independent: the silent orbit of a state may be infinite while the reachable states are finite, or
  conversely. Either failure is genuine non-termination of the construction.
\end{remark}

\begin{remark}[Uniqueness does not characterise halting]\label{rem:notmunu}
  One might hope the tabling halts exactly when the least and greatest fixpoint coincide. They always
  coincide (Theorem~\ref{thm:unique} gives a unique fixpoint for \emph{every} term), yet the tabling does
  not halt on every term, so $\mu=\nu$ cannot be the criterion. The witness is the stream of all
  naturals, $\mathit{nats} = Y(\lambda s.\lambda n.\,\mathtt{cons}\,n\,(s\,(\mathtt{succ}\,n)))\,\underline{0}$:
  it is guarded, so its behaviour is that unique fixpoint, a well-defined infinite stream, but the stream
  is \emph{non-rational} (its tails are pairwise distinct), so it has no finite representation and the
  tabling diverges. Halting is governed by rationality, which is strictly stronger than uniqueness: the
  rational behaviours are a proper subclass of the (always unique) behaviours.
\end{remark}

\begin{theorem}[Finite-representation barrier; optimality]\label{thm:optimal}
  By a \emph{finite representation} of a behaviour $t \in \nu\Beh$ we mean a finite $\Sigma$-graph whose
  unfolding is $t$ (a first-order datum, to be distinguished from a finite \emph{generator}, a term that
  reduces to $t$). The statement is for the deterministic skeleton; for a branching $M$ the labels carry
  the $M$-structure and $\rfp\Beh$ is the corresponding rational fixpoint, and the argument lifts
  unchanged. Then:
  \begin{enumerate}
    \item \emph{(Barrier.)} $t$ has a finite representation iff $t$ is rational (Definition~\ref{def:rational},
      i.e.\ $t \in \rfp\Beh$); hence no procedure halts with a finite representation of a non-rational
      behaviour, even though every behaviour has a finite generator.
    \item \emph{(Soundness.)} Whenever the tabling halts it outputs a finite representation
      (Theorems~\ref{thm:termination}, \ref{thm:adequacy}), so it halts only on rational behaviours.
    \item \emph{(Attainment.)} The tabling halts iff finitely many positions are reached up to the identity
      it folds on. Under structural identity the halting set is a proper subclass of the rational
      behaviours, the gap being the dead-subterm witnesses ($Y\,F\,\underline{0}$ with
      $F=\lambda s.\lambda n.\,\mathtt{cons}\,0\,(s\,(\mathtt{succ}\,n))$); folding additionally any
      positions of equal behaviour enlarges it, and in the limit, folding all behaviourally equal
      positions (a coarser identification than structural identity, and in the $\lambda$-instance not
      decidable, so an ideal limit rather than the relation the implementation folds on), the tabling
      halts iff the behaviour is rational. By (1) no sound identity pushes the halting set beyond the
      rational fragment.
  \end{enumerate}
\end{theorem}
\begin{proof}
  \emph{(1)} A finite representation is a finite $\Sigma$-graph; its unfolding has at most one distinct
  subtree per node, hence finitely many, so it is rational. Conversely a rational $t$, having finitely
  many distinct subtrees, is the unfolding of the finite graph on those subtrees with an edge to each
  immediate subtree of a root; this is the standard equivalence for regular trees \citeTODO{courcelle},
  and the rational behaviours are exactly $\rfp\Beh$ (Definition~\ref{def:behfun}). Hence a non-rational
  $t$ is the unfolding of no finite $\Sigma$-graph, so no procedure can output one. The generator clause
  is immediate: the term $M$ is a finite generator of $\behav(M)$ for every $M$, rational or not. The
  barrier is against any \emph{regular} (first-order, finitely presented) representation, each of which
  unfolds to only rational trees; only a Turing-complete generator, such as the term itself, escapes it.

  \emph{(2)} If the tabling halts then by Theorem~\ref{thm:termination} the reached positions are finite
  and the result is a finite graph in $\rfp\Beh$, which by Theorem~\ref{thm:adequacy} unfolds to
  $\behav(M)$; so a halt witnesses $\behav(M) \in \rfp\Beh$.

  \emph{(3)} The tabling descends through positions up to the folding identity and tables them, so it
  halts iff those classes are finite (Theorem~\ref{thm:termination}). Structural identity can keep
  infinitely many classes while the behaviour is rational: in the cited witness every tail
  $Y\,F\,\underline{k}$ is a distinct closed term, yet all denote the one constant-$0$ stream, which is
  rational; so the halting set is a proper subclass of the rational behaviours. An identity that only
  ever identifies positions of equal behaviour is sound (it changes no tree, Theorem~\ref{thm:adequacy})
  and yields fewer classes, hence a larger halting set; for the coarsest such identity two positions are
  identified exactly when their subtrees are equal, so the classes are the distinct subtrees of
  $\behav(M)$, finite iff $\behav(M)$ is rational. That identity therefore halts iff rational, and by (1)
  no sound identity can do better. That coarsest identity is effective on the rational fragment, where
  each fold is decided by cycle-detecting bisimulation; the implementation here ships structural identity
  instead, which sits strictly below it, the gap being the dead-subterm witnesses.
\end{proof}

\begin{remark}[The optimum is relative to the denotation]\label{rem:fixeddenotation}
  The optimum in Theorem~\ref{thm:optimal} is among procedures that \emph{preserve the denotation}, the
  tree. A fold changes no tree exactly when it identifies only behaviourally equal positions, so
  behavioural equivalence is the coarsest denotation-preserving identity, and it is what attains the
  rational ceiling (structural identity, fixed by the syntax, sits below it; behavioural equivalence,
  relative to the chosen normalisation, sits at it). Folding more would identify positions of
  \emph{different} behaviour: that is a coarser normalisation, a different tree, hence a different
  denotation. One can then terminate on more, even on $\mathit{nats}$ by collapsing its distinct tails,
  but one is computing a different semantics. Thus Theorem~\ref{thm:optimal}(1) is the barrier \emph{at a
  fixed denotation}; termination beyond it is bought only by changing what is computed.
\end{remark}

\begin{remark}[Optimal semi-decision]\label{rem:semidec}
  Rationality is undecidable: given a machine $T$, the term that emits the constant-$0$ stream while
  simulating $T$ and switches to emitting $0,1,2,\dots$ once $T$ halts has a rational behaviour iff $T$
  never halts, so rationality is co-r.e.-hard. Hence the terminating inputs cannot be pre-selected. The tabling is thus a \emph{sound semi-decision procedure}
  for rationality: on the rational fragment (up to its identity, Theorem~\ref{thm:optimal}(3)) it halts
  and returns the finite representation, and elsewhere it diverges, which by Theorem~\ref{thm:optimal}(1)
  is the most any procedure can achieve.
\end{remark}

\begin{remark}[Determinism: reduction strategies]\label{rem:determinism}
  When $M = \Delaym$, an equation morphism is a deterministic observe-or-step system
  $X \to F X + X$ (one silent successor, divergence the only effect), the behaviour functor is $\Beh = F$
  with $\nu\Beh$ the ordinary $F$-trees, and the silent collapse decides an unproductive cycle as the
  single value $\bot$. This is the case of a \emph{deterministic reduction strategy}: head and
  weak-head normalisation (Section~\ref{sec:bridge}), and equally call-by-name and call-by-value head
  reduction, differ only in the choice of $F$ and of which redex the step contracts, never in the
  $M = \Delaym$ shape. Richer $M$ (non-deterministic, probabilistic, stateful) leaves the architecture
  intact but replaces the cycle-to-$\bot$ collapse with the corresponding $M$-least fixpoint, for
  $M = \Pow$ a set-valued one.
\end{remark}

% =====================================================================
\section{The behaviour has a unique fixpoint: least equals greatest}\label{sec:unique}
% =====================================================================
The behaviour was introduced coinductively, as the unfolding of $\outm$ into $\nu\Beh$
(Definition~\ref{def:behreadout}). It is equally the order-theoretic least fixpoint of a behaviour
functional, and the two coincide: that functional is \emph{guarded}, so it has a \emph{unique}
fixpoint, and no least-versus-greatest choice arises. Once $\outm$ has exposed a constructor,
induction and coinduction agree.

\begin{definition}[Behaviour functional]\label{def:readoutfun}
  The \emph{behaviour functional} $\Psi\colon (X \to D) \to (X \to D)$ reads one $\Beh$-layer off the
  structure map and defers to its argument for the children:
  \[
    \Psi(g)(n) =
    \begin{cases}
      \bot & \outm(n) = \bot,\\
      \sigma\big(g(n_1),\dots,g(n_k)\big) & \outm(n) \text{ exposes the layer } \sigma(n_1,\dots,n_k),
    \end{cases}
  \]
  the $M$-structure of the exposed layer (for a branching $M$) carried at the root, above the children.
  By Proposition~\ref{prop:cpo} the pointwise order on $X \to D$ is a pointed $\omega$-CPO and $\Psi$ is
  monotone and $\omega$-continuous, so $\operatorname{lfp}(\Psi) = \bigsqcup_{k} \Psi^{k}(\lambda n.\,\bot_D)$
  exists.
\end{definition}

\begin{definition}[Tree metric]\label{def:metric}
  For $t,u \in D$ put $\mathrm{d}(t,u) = 0$ if $t = u$, and otherwise $\mathrm{d}(t,u) = 2^{-\ell}$ where
  $\ell$ is the least length $|p|$ of a position $p \in \mathrm{dom}(t) \cup \mathrm{dom}(u)$ at which $t$
  and $u$ differ ($p$ lies in only one domain, or $t(p) \neq u(p)$). Lift it to $X \to D$ by
  $\mathrm{d}(g,h) = \sup_{n} \mathrm{d}(g(n),h(n))$.
\end{definition}

\begin{proposition}[The tree space is complete]\label{prop:metric}
  $(D,\mathrm{d})$ is a complete ultrametric space, and hence so is $(X \to D, \mathrm{d})$ under the
  supremum metric \citeTODO{arnold-nivat}. A Cauchy sequence of trees stabilises at every fixed depth,
  so it has a level-wise limit in $D$; the metric topology agrees with the $\omega$-CPO of
  Proposition~\ref{prop:cpo} in that the $\mathrm{d}$-limit of a $\sqsubseteq$-increasing sequence is its
  least upper bound.
\end{proposition}

\begin{lemma}[The behaviour functional is guarded]\label{lem:guarded}
  $\Psi$ is a contraction with factor $\tfrac{1}{2}$:
  $\mathrm{d}(\Psi(g),\Psi(h)) \le \tfrac{1}{2}\,\mathrm{d}(g,h)$.
\end{lemma}
\begin{proof}
  Fix a state $n$. The root label of $\Psi(g)(n)$ is determined by $\outm(n)$ alone, independently of
  $g$: it is $\bot$ when $\outm(n) = \bot$, and the operation $\sigma$ when $\outm(n)$ exposes
  $\sigma(n_1,\dots,n_k)$. So $\Psi(g)(n)$ and $\Psi(h)(n)$ always share the root ($|p| = 0$) label and
  can differ only strictly below it. If $\outm(n) = \bot$ or a nullary layer, they are equal and
  $\mathrm{d}(\Psi(g)(n),\Psi(h)(n)) = 0$. Otherwise a shortest position at which they differ is one
  descent step into some child $n_i$ followed by a shortest differing position of $g(n_i)$ and $h(n_i)$,
  so its length is $1$ greater, giving
  \[
    \mathrm{d}\big(\Psi(g)(n),\Psi(h)(n)\big) = \tfrac{1}{2}\,\max_{i}\,\mathrm{d}\big(g(n_i),h(n_i)\big)
    \le \tfrac{1}{2}\,\mathrm{d}(g,h).
  \]
  Taking the supremum over $n$ gives $\mathrm{d}(\Psi(g),\Psi(h)) \le \tfrac{1}{2}\,\mathrm{d}(g,h)$.
\end{proof}

\begin{theorem}[Least equals greatest]\label{thm:unique}
  $\Psi$ has a unique fixpoint, and it is simultaneously
  \begin{enumerate}
    \item the order-theoretic least fixpoint $\operatorname{lfp}(\Psi) = \bigsqcup_{k} \Psi^{k}(\lambda n.\,\bot_D)$, and
    \item the coinductive unfolding $\behav$, the unique $\Beh$-coalgebra morphism $X \to \nu\Beh$ over
      $\outm$ (Definition~\ref{def:behreadout}).
  \end{enumerate}
  Hence the least-fixpoint reading and the coinductive reading denote the same object: there is no
  least-versus-greatest choice to make.
\end{theorem}
\begin{proof}
  By Lemma~\ref{lem:guarded}, $\Psi$ is a contraction on the complete metric space $(X \to D, \mathrm{d})$
  (Proposition~\ref{prop:metric}); by the Banach fixed-point theorem \citeTODO{banach} it has a unique
  fixpoint $r$, and $\Psi^{k}(g_0) \to r$ from every starting point $g_0$.

  \emph{(1)} Take $g_0 = \lambda n.\,\bot_D$. The iterates $\Psi^{k}(g_0)$ are $\sqsubseteq$-increasing
  ($\Psi$ is monotone and $g_0$ is least) and converge to $r$ in $\mathrm{d}$; the $\mathrm{d}$-limit of a
  $\sqsubseteq$-increasing sequence is its least upper bound (Proposition~\ref{prop:metric}), so
  $\bigsqcup_{k} \Psi^{k}(g_0) = r$, which by $\omega$-continuity is $\operatorname{lfp}(\Psi)$. Thus
  $\operatorname{lfp}(\Psi) = r$.

  \emph{(2)} A function $g$ is a fixpoint of $\Psi$ iff $g(n) = \sigma(g(n_1),\dots,g(n_k))$ whenever
  $\outm(n)$ exposes $\sigma(n_1,\dots,n_k)$ and $g(n) = \bot$ when $\outm(n) = \bot$, which is exactly
  the equation defining the $\Beh$-coalgebra morphism $X \to \nu\Beh$ over $\outm$. By finality that
  morphism $\behav$ exists and is the unique such; being a fixpoint of $\Psi$, it equals $r$.

  Therefore $\operatorname{lfp}(\Psi) = r = \behav$.
\end{proof}

\begin{remark}[Algebraic compactness]\label{rem:compact}
  Theorem~\ref{thm:unique} is the operational face of a categorical fact. Over a \emph{discrete} carrier
  the initial $\Beh$-algebra and the final $\Beh$-coalgebra differ, $\mu\Beh \neq \nu\Beh$ (for the
  stream functor, the finite lists versus the finite-or-infinite ones), so there least and greatest
  fixpoint genuinely diverge. The behaviour functional instead lives over the $\bot$-completed tree domain $D$, where a
  locally continuous functor is \emph{algebraically compact}, $\mu\Beh \cong \nu\Beh$
  \citeTODO{smyth-plotkin, freyd}, and the unique fixpoint of Theorem~\ref{thm:unique} is that one
  common object. So ``least equals greatest'' is a property of the $\bot$-completion and the guarded
  behaviour functional, \emph{not} of the discrete signature: completing with $\bot$ and reading guardedly is exactly
  what collapses induction and coinduction.
\end{remark}

\begin{remark}[The parameter is the approximation order; the result is about computable ADTs]\label{rem:metalanguage}
  What forces the unique fixpoint of Theorem~\ref{thm:unique} is the \emph{approximation order} on terms
  (Definition~\ref{def:order}, equivalently its metric, Definition~\ref{def:metric}), under which the
  guarded behaviour functional is a contraction; the iteration plays no role. This order is intrinsic to the
  \emph{signature}, not to the $\lambda$-calculus: the Scott encoding supplies it automatically at the
  implementation level (a constructor is observed, a divergence is $\bot$), but it can be defined
  directly, with no reference to $\lambda$-terms, as the ideal completion of the finite partial
  $\Sigma$-trees. Fixing it determines at once the order and the unique fixpoint of \emph{every} data
  structure the calculus describes: a structure presented by a $\lambda$-term \emph{is} the tree read off
  $\outm$, so it inherits both. Theorem~\ref{thm:unique} is therefore not really about the
  $\lambda$-calculus. For any finitary signature $F$, every computable abstract data type, taken with its
  approximation order, has a unique fixpoint, its initial algebra and final coalgebra coinciding
  (Remark~\ref{rem:compact}); the $\lambda$-calculus is one realization of that order, not its source.
  One never designates an iteration operator or selects a least solution \emph{for the data}; the least
  fixpoint survives only one level down, in the evaluation, where the silent collapse gives a divergent
  computation its forced value $\bot$ (Remark~\ref{rem:leastessential}).
\end{remark}

\begin{remark}[Where the least fixpoint is indispensable]\label{rem:leastessential}
  Uniqueness is a property of the behaviour functional, after $\outm$ has exposed a constructor. The structure map
  $\outm = \iter{\stepm}$ itself is \emph{unguarded}: a silent cycle returns with no constructor exposed,
  so its defining operator (Definition~\ref{def:out}) is not contractive, and there least genuinely
  differs from greatest. The least choice is the correct and the only computable one: $\Omega$ has
  $\outm(\Omega) = \bot$ (least), whereas a greatest reading would assign it an arbitrary non-$\bot$
  shape. So the construction is least-fixpoint throughout, and the single place the choice does real work
  is the collapse of divergence to $\bot$; once a layer is exposed, the behaviour functional is guarded and least
  equals greatest.
\end{remark}

% =====================================================================
\section{Bridge: head normalisation is an equation morphism}\label{sec:bridge}
% =====================================================================
Head normalisation of the pure $\lambda$-calculus is the instance $M = \Delaym$ (pure divergence,
Example~\ref{ex:effects}) with $F$ the $\lambda$-shape signature. It is the instance that bears out the
opening claim, so we give it in full; everything specialises the general development.

\begin{definition}[The $\lambda$-shape signature]\label{def:lamsig}
  States are pure $\lambda$-terms in de Bruijn form (no $\alpha$-renaming), interned so that
  structurally equal terms are one state. The signature is
  \[
    F X = \big\{\Varsh_i\big\}_{i\in\mathbb{N}} \;+\; \big\{\Lamsh(x)\big\}_{x\in X} \;+\; \big\{\Appsh(x,x')\big\}_{x,x'\in X},
  \]
  a nullary $\Varsh_i$ per de Bruijn index, a unary $\Lamsh$, a binary $\Appsh$. There is \emph{no}
  recursion binder; the repetition that tabling folds is produced by ordinary terms such as the
  $Y$ combinator.
\end{definition}

\begin{definition}[Weak-head reduction as the equation morphism]\label{def:shbridge}
  The equation morphism $\stepm \colon X \to \Delaym(F X + X)$ is weak-head reduction:
  \[
    \stepm(M) =
    \begin{cases}
      \mathrm{now}\,\mathrm{inl}(\Varsh_i) & M = \Varsh_i\\
      \mathrm{now}\,\mathrm{inl}(\Lamsh(N)) & M = \Lamsh(\lambda x.N)\\
      \mathrm{now}\,\mathrm{inr}\big(\mathit{body}[x\mathbin{:=}a]\big) & M = \Appsh(f,a),\ f \text{ exposes } \Lamsh(\lambda x.\mathit{body})\\
      \mathrm{now}\,\mathrm{inl}(\Appsh(f,a)) & M = \Appsh(f,a),\ f \text{ exposes a rigid head.}
    \end{cases}
  \]
  A head redex is a silent step ($\mathrm{inr}$, contract and continue); a variable, an abstraction,
  or an application with rigid head exposes a layer ($\mathrm{inl}$). The head $f$ is inspected through
  its own structure map $\outm(f)$, so the clauses are mutually recursive through $\outm$, exactly the
  silent iteration of Definition~\ref{def:out}.
\end{definition}

\begin{proposition}[The weak-head structure map is the weak-head shape]\label{prop:shisout}
  For this $\stepm$, the structure map $\outm = \iter{\stepm}$ (Definition~\ref{def:out}) is the
  weak-head shape relation $\Sh$: $\outm(M)$ is the outermost constructor of $M$ after weak head
  reduction, and $\outm(M) = \bot$ exactly when weak head reduction of $M$ does not terminate. The
  behaviour $\behav$ (Definition~\ref{def:behreadout}) is the L\'evy--Longo tree. Taking the structure
  map to be head reduction \emph{under} $\lambda$ instead, a ``$\lambda$'' over a body with no head
  normal form being itself $\bot$, gives the B\"ohm tree (Remark~\ref{rem:boehm}). Proof: specialise
  Propositions~\ref{prop:outwd}, \ref{prop:cpo} and Theorem~\ref{thm:adequacy} to $M = \Delaym$.
\end{proposition}

\begin{remark}[Both readings, same convergence]\label{rem:boehm}
  The weak-head and head structure maps are two instances over the one carrier, the one tabling, and
  the one behaviour map, differing only in the $\lambda$ clause. They terminate on exactly the same terms:
  head reduction performs eagerly, along the head spine, the very weak-head reductions the L\'evy--Longo
  behaviour computation performs lazily down the $\lambda$-spine, so the reachable nodes and the reductions coincide.
  They differ only denotationally, at a $\lambda$ whose body has no head normal form: L\'evy--Longo
  keeps it (reading $\lambda x.\,\Omega$ as $\lambda x.\,\bot$), B\"ohm bottoms the whole term (reading
  it as $\bot$). The choice of normalisation is therefore denotational, not a difference in
  convergence. The accompanying implementation\anon[~(supplementary
  material)]{~\cite{colambda}} realises both as least-fixpoint structure maps over the one
  carrier and the one behaviour map, and checks both facts: $\lambda x.\,\Omega$ reads out as
  $\lambda x.\,\bot$ under weak head and as $\bot$ under head, while both readings terminate on the same
  terms and diverge on the same non-rational stream.
\end{remark}

\begin{example}[Guarded: the cyclic stream $Y\,(\mathtt{cons}\;0)$]\label{ex:cons}
  Write $r = \mathtt{cons}\;0\;r$ as $Y\,(\mathtt{cons}\;0)$, with the Scott encoding
  $\mathtt{cons} = \lambda h.\lambda t.\lambda c.\lambda n.\,c\,h\,t$. Then $\outm(r)$ exposes
  $\lambda c.\lambda n.\,c\,0\,r'$ in finitely many silent steps without re-requiring $\outm(r)$: the
  recursive occurrence enters only as data. The repeating cell is one interned state, so tabling
  folds at the guarded re-entry,
  \[
    \rho \;=\; \lambda c.\,\lambda n.\,(c\cdot 0)\cdot\rho,
  \]
  the finite representation of the rational L\'evy--Longo tree
  $\lambda c.\lambda n.\,(c\,0)\,(\lambda c.\lambda n.\,(c\,0)\,(\cdots))$, a node pointing to itself,
  where head reduction unfolds forever. No recursion binder is needed; $Y$ produces the repetition that
  tabling folds.
\end{example}

\begin{example}[Unproductive: $\Omega$ and $Y\,(\lambda x.x)$]\label{ex:xx}
  $\Omega = (\lambda x.x\,x)(\lambda x.x\,x)$ steps to itself with no constructor exposed, returning to
  the same interned state; the silent collapse $\outm(\Omega)$ has nothing to expose, so $\outm(\Omega)
  = \bot$ and $\Omega$ reads out as $\bot$, decided in finite time where reduction loops. The same holds
  for $Y\,(\lambda x.x)$, the pure-$\lambda$ rendering of $\mathtt{letrec}\;x = x$.
\end{example}

\begin{remark}[Where this gains over call-by-need]\label{rem:callbyneed}
Lazy (call-by-need) evaluation memoises only thunks that are aliased in the heap, so it folds a
recurrence only when the recurring subterm is physically the same cell. On a term with a finite normal
form it then agrees with the tabling on the value and differs only in cost: the tabling identifies
subcomputations that are structurally equal yet separately built, collapsing for instance the
exponential edit-distance recursion to its $O(mn)$ distinct suffix pairs, the dynamic-programming
table, with no memoisation written. The gain in convergence is elsewhere. On a behaviour that is
infinite yet reaches finitely many states, such as $Y\,(\mathtt{cons}\;0)$ (Example~\ref{ex:cons}) or
$\mathtt{map}\,\mathtt{succ}\,(Y\,(\mathtt{cons}\;0))$, whose rebuilt spine shares no cell with its
source, call-by-need diverges while the tabling returns the finite cyclic graph; and on an
unproductive term such as $\Omega$ (Example~\ref{ex:xx}) the tabling decides $\bot$ in finite time
where call-by-need loops. This gain is exactly a decidable fragment of the rational ceiling of
Theorem~\ref{thm:optimal}: the tabling on structural identity halts on the behaviours that reach
finitely many states, and the residual gap to the full rational ceiling is undecidable
(Remark~\ref{rem:semidec}).
\end{remark}

% =====================================================================
\section{Applications}\label{sec:application}
% =====================================================================
\subsection{Decidability and dissolved problems}\label{sec:variants}
On the rational fragment the behaviour is a finite graph, so its equality is decidable, and a cluster
of problems that need added machinery in other settings does not arise.

\begin{corollary}[Decidable behaviour equality on the rational fragment]\label{cor:decidable}
  Restricted to states from which finitely many states are reachable (Theorem~\ref{thm:termination}),
  equality of behaviours $\behav(x) = \behav(y)$ is decidable. Each is presented by a finite cyclic
  graph; the graphs are computable, since the silent collapse terminates at each reached state
  (Proposition~\ref{prop:outwd}) and the reachable states close after finitely many steps, while state
  identity is decidable (structural equality of states). Equality of two rational graphs is then decided
  by the standard cycle-detecting bisimulation \citeTODO{courcelle}.
\end{corollary}

\begin{corollary}[Hard problems dissolved]\label{cor:dissolved}
  Several problems that need added machinery elsewhere do not arise here. Observing sharing or cyclicity
  in a pure language needs an observable-sharing primitive or a circular-program idiom
  \citeTODO{gill, bird}; here sharing is state identity. Representing a cyclic structure needs an
  explicit recursion binder or a cyclic $\lambda$-calculus \citeTODO{ariola-klop}; here the equation
  morphism generates the repetition and tabling folds it, with no binder added
  (Section~\ref{sec:bridge}). Deciding equality over cyclic structures needs a separate unification
  discipline or recursive-type system \citeTODO{colmerauer, amadio-cardelli, brandt-henglein}; here it
  is cycle detection over states (Corollary~\ref{cor:decidable}). Terminating a recursive query over a
  cyclic datum needs a tabled evaluator \citeTODO{tamaki-sato, tabling-slg}; here the query's own
  recurrences fall into the same states and fold (Section~\ref{sec:reach}).
\end{corollary}

\subsection{Structural computation over cyclic data}\label{sec:reach}
Constructing the cyclic datum is already the new step. The pure term $Y\,(\mathtt{cons}\;0)$ denotes
the infinite stream $\mathtt{cons}\;0\;(\mathtt{cons}\;0\;\cdots)$, with no finite form under head
reduction; up to identity it has finitely many states, so tabling $\outm$ gives it a finite
\emph{circular} representation, a cell whose tail points back to itself (Example~\ref{ex:cons}). A
circular linked list from a pure $\lambda$-term is what the standard semantics cannot build.

A recursive query over that datum is an ordinary $\lambda$-function, and the guarded/unproductive
split of Definition~\ref{def:guarded} governs the \emph{query} as it governs the data. The na\"ive
walk
\[
  \mathtt{reach} \;=\; Y\,\big(\lambda\mathit{self}.\,\lambda\mathit{lst}.\;
  \mathit{lst}\;(\lambda h.\lambda t.\,\mathit{self}\;t)\;\present\big)
\]
recurses down a Scott stream looking for $\mathtt{nil}$. Applied to $r = \mathtt{cons}\;0\;r$, each
recursive call builds $\mathtt{reach}\;t$; because $r$ has finitely many states, these fall into
finitely many states, so tabling folds the recursion. Head reduction, which remembers nothing,
builds a fresh application at every step and diverges; tabling folds at the repeated state and
returns a definite value (here $\bot$: no $\mathtt{nil}$ is reached, so the walk is itself an
unproductive cycle), in finite time. It is memoised for free, since identical states share their
subcomputations.

\paragraph{Mapping a cyclic list}
The guarded arm produces rather than decides. The textbook
$\mathtt{map}\,f = Y\,(\lambda\mathit{self}.\,\lambda\mathit{lst}.\;\mathit{lst}\;
(\lambda h.\lambda t.\,\mathtt{cons}\;(f\,h)\;(\mathit{self}\;t))\;\mathtt{nil})$ has nothing
cycle-aware in it. Its recursive call sits under a $\mathtt{cons}$, so it is guarded: applied to
$r = \mathtt{cons}\;0\;r$, the recursive application $\mathtt{map}\,f\;t$ re-enters the same state
$\mathtt{map}\,f\;r$, and tabling folds it. The result is a finite circular list
$\rho = \mathtt{cons}\;(f\,0)\;\rho$, every element transformed, where head reduction would map the
infinite stream forever. The fold is lazy knot-tying: each cell is registered before its tail is
recursed into, so the table supplies the back edge. On a finite list the same $\mathtt{map}$ is the
ordinary one. The implementation\anon[~(supplementary material)]{~\cite{colambda}} reproduces
both arms.

\subsection{Memoising arbitrary programs}\label{sec:memo}
More generally, $\behav$ is the denotation of a \emph{memoised} run of any program over the states.
Memoisation needs a comparable key on which to recognise a repeated call; first-order states supply
one, higher-order values do not (value equality of functions is undecidable). The tabling that computes the behaviour
is that memoisation: it is what lets a re-entry be recognised at all. Head reduction, which remembers
nothing, unfolds a cyclic state forever; tabling $\outm$ recognises the re-entry and folds. The two
are the unfolding and its limit, separated by how much they remember.

% =====================================================================
\section{Related work}\label{sec:related}
% =====================================================================
\paragraph{Coalgebra, final semantics, and rational fixpoints.}
Reading behaviour as the unique map into a final coalgebra is the coalgebraic method
\citeTODO{rutten, jacobs}; the finite cyclic graphs are the rational fixpoint between the initial
algebra and the final coalgebra \citeTODO{adamek-milius-velebil-rational}. Partition refinement
\citeTODO{hopcroft, paige-tarjan}, and its generic coalgebraic form CoPaR \citeTODO{copar,
wissmann-modular-partition-refinement} for any zippable functor on $\mathbf{Set}$, solves a different
problem: it is \emph{given} a finite coalgebra and \emph{minimises} it, quotienting its states by
behavioural equivalence (the coarsest bisimulation, the kernel of the map into the final coalgebra),
a comparison among the states of an already-given system. We do not minimise or compare two systems; we
\emph{solve} one: a $\lambda$-term is read as a coalgebraic equation, and we compute its solution, and on
the rational fragment a finite representation of that one solution. The only equality we use is
\emph{structural identity} of nodes \emph{within that single solution}, and only to recognise a
re-entrant cycle and tie a back edge (decidable because the nodes are first-order,
Section~\ref{sec:intro}); we never quotient by behavioural equivalence, which on $\lambda$-terms would be
B\"ohm-tree equality and is undecidable. The two tasks are orthogonal and could compose: CoPaR could
minimise the finite graph we produce. Our behaviour, finally, is the first-order $\Sigma$-tree a state
unfolds to, not the extensional behaviour of a higher-order value.

\paragraph{Iteration, the delay monad, and Elgot monads.}
Iteration in a (complete) Elgot monad
\citeTODO{adamek-milius-velebil, goncharov-elgot} solves guarded recursion abstractly, and for
$M = \Delaym$ the delay/partiality monad \citeTODO{capretta} makes divergence a value; but these
accounts fix the value categorically and neither compute it nor give it a finite representation on
the rational fragment. Generalised resumption monads $M(F(-) + -)$ \citeTODO{piróg-gibbons} supply the
type of such systems but no procedure over them.

\paragraph{Least Herbrand models and tabling.}
The least fixpoint of a monotone consequence operator is the least Herbrand
model of a definite specification \citeTODO{vanemden-kowalski}; plain resolution diverges on a cyclic
query, whereas a tabled evaluator \citeTODO{tamaki-sato, tabling-slg} computes the least fixpoint by
completing unproductive cycles instead of looping. Tabling has also been extended to store rational
(cyclic) terms under a canonical representation and to drive coinductive logic programming
\citeTODO{mantadelis-rocha-moura}, but there the cyclic term is already produced by occurs-check-free
unification. In all of these, tabling stays within logic programming and is not carried to an arbitrary
signature functor.

\paragraph{Rational trees and recursive types.}
Deciding structural queries over cyclic structures by cycle detection is long established: rational
(regular) trees and their algebra \citeTODO{courcelle} underlie cyclic unification
\citeTODO{colmerauer}, and equirecursive recursive types give terminating subtyping and equality by
detecting cycles in a regular tree \citeTODO{amadio-cardelli, brandt-henglein}. Each, however, handles
only the guarded-cycle case and is a separate discipline tied to one domain, unification or subtyping,
not a single uniform construction.

\paragraph{$\lambda$-trees and infinitary $\lambda$-calculus.}
The Berarducci, L\'evy--Longo, and B\"ohm trees form a hierarchy by their meaningless sets
\citeTODO{berarducci, levy-longo, barendregt}, situated uniformly by the infinitary
$\lambda$-calculus \citeTODO{kennaway-klop-sleep-devries}; they are studied as infinite normal forms,
with no finite representation of the rational ones. The Berarducci tree retains $\bot$-headed
applications and is non-monotone, so it admits no least-fixpoint presentation. The lazy
$\lambda$-calculus \citeTODO{abramsky-ong} observes weak head normal forms but does not fold the trees
it produces.

\paragraph{Cyclic $\lambda$-terms, sharing, and recursion.}
Graph reduction \citeTODO{wadsworth, turner-combinators} evaluates $\lambda$-terms over (possibly
cyclic) graphs, and term graph rewriting \citeTODO{barendregt-term-graph, kennaway-graph-adequacy}
relates finite cyclic graphs to the rational trees they unfold to; but a cycle is tied only where
recursion is primitive, a \texttt{letrec} or a built-in fixpoint, whereas a pure fixpoint combinator
has its body copied at each $\beta$-step and unfolds into an unbounded spine. Cyclic graphs / explicit
recursion \citeTODO{ariola-klop} and call-by-need \citeTODO{ariola-felleisen} likewise provide sharing
only as an added construct, and a pure functional language cannot even observe sharing or cyclicity
without an addition such as observable sharing \citeTODO{gill} or a circular-program idiom
\citeTODO{bird}. Across all of these, an ordinary recursive function over a cyclic list, such as a
$\mathtt{map}$, reduces to an unbounded spine rather than to a cycle.

\paragraph{Higher-order model checking and defunctionalization.}
Deciding properties of trees generated by higher-order programs \citeTODO{ong-hors, kobayashi}
answers queries about $\lambda$-generated infinite trees, but proceeds by automata and games over
B\"ohm-tree-like denotations and produces no finite representation of the tree itself.
Defunctionalization \citeTODO{reynolds} compiles a higher-order program into a first-order
interpreter, a syntactic translation that yields no semantic \emph{object}.

\paragraph{A common limitation.}
Terminating structural computation over cyclic data by cycle detection is decades old in logic
programming \citeTODO{colmerauer, tamaki-sato} and recursive-type theory \citeTODO{amadio-cardelli,
brandt-henglein}, but in every case it is confined to one domain and applies only to data already
presented in cyclic form; none reads it as the behaviour of an arbitrary system, and none folds a pure
term that merely generates the cycle.

% =====================================================================
\section{Conclusion}\label{sec:conclusion}
% =====================================================================
The behaviour of an equation morphism is read by tabling a one-step structure map, the structure map
itself the least fixpoint of the silent iteration. The unfolding diverges on a cyclic state; tabling
that map on first-order state identity folds a guarded cycle into a finite graph and decides an
unproductive cycle as $\bot$, both in finite time on the rational fragment. Head normalisation of
the pure $\lambda$-calculus is one instance, the one that shows a pure term can carry and compute over
cyclic data with nothing added to the calculus. The two parameters, the effect monad of the internal
process and the signature of the exposed layer, range over deterministic, non-deterministic,
probabilistic, and stateful processes alike.

% =====================================================================
\appendix
\section{The behaviour domain (Proposition~\ref{prop:cpo})}\label{app:domain}
\begin{proof}[Proof of Proposition~\ref{prop:cpo}]
  \emph{Least element.} For any $t \in D$, $\mathrm{dom}(\bot_D) = \{\varepsilon\} \subseteq
  \mathrm{dom}(t)$ and $\bot_D(\varepsilon) = \bot$, so $\bot_D \sqsubseteq t$ by
  Definition~\ref{def:order}.

  \emph{Directed completeness.} Let $\{t_i\}_{i\in I}$ be $\sqsubseteq$-directed. Define $t$ by
  $\mathrm{dom}(t) = \bigcup_i \mathrm{dom}(t_i)$ and $t(p) = t_i(p)$ for any $i$ with
  $t_i(p) \neq \bot$, else $t(p) = \bot$. This is well-defined: if $t_i(p) \neq \bot \neq t_j(p)$,
  directedness gives $t_\ell \sqsupseteq t_i, t_j$, and a non-$\bot$ label is preserved upward
  (Definition~\ref{def:order}), so $t_i(p) = t_\ell(p) = t_j(p)$. And $t \in D$: $\mathrm{dom}(t)$ is
  prefix-closed as a union of prefix-closed sets, and arity-respecting because a non-$\bot$ label at $p$
  agrees with some $t_i$, whose children of $p$ are exactly those the label prescribes, and $\bot$ has
  no children. Finally $t = \bigsqcup_i t_i$: each $t_i \sqsubseteq t$, and any upper bound $u$ has
  $t \sqsubseteq u$, since at $p$ with $t(p) \neq \bot$ the witnessing $t_i$ gives
  $t(p) = t_i(p) = u(p)$. The $M$-structure at each node is completed $M$-pointwise, where $M A$ is a
  pointed $\omega$-CPO by Definition~\ref{def:elgot}. In particular every $\omega$-chain has a least
  upper bound.
\end{proof}

\section{The structure map (Proposition~\ref{prop:outwd})}\label{app:out}
\begin{proof}[Proof of Proposition~\ref{prop:outwd}]
  Write $\Phi$ for the operator of Definition~\ref{def:out}, on $X \to M(F X)$ ordered pointwise over
  the $\omega$-CPO structure of $M(F X)$ (Definition~\ref{def:elgot}).

  \emph{Monotone and continuous.} $\Phi(g)(x) = \big([\,\eta\circ\mathrm{inl},\,g\,]^{*}\big)(\stepm(x))$
  is built from $\stepm(x)$ (a constant in $g$) by Kleisli extension of $[\eta\circ\mathrm{inl}, g]$,
  which depends on $g$ only through the $\mathrm{inr}$ branch and is $\omega$-continuous in $g$ because
  $(-)^{*}$ and copairing are $\omega$-continuous (Definition~\ref{def:elgot}). Hence $\Phi$ is
  $\omega$-continuous, and by Kleene $\outm = \bigsqcup_k \Phi^k(\lambda x.\,\bot)$ is its least
  fixpoint.

  \emph{Characterisation.} The iterates resolve layers by running silent steps: $\Phi^k(\bot)(x) \neq
  \bot$ iff the silent process from $x$ exposes an $\mathrm{inl}$ layer within $k$ steps. \TODO{spell out
    the step measure: an induction on $k$ matching ``$\Phi^k(\bot)(x) \neq \bot$'' with ``$x$ exposes a
    layer within $k$ silent steps'', concluding $\outm(x) = \bot$ iff the silent process from $x$ exposes
  no layer. The statement and both endpoints are fixed; only the bookkeeping of the measure remains.}
  For $M = \Delaym$ this says $\outm(x) = \bot$ exactly when weak head reduction does not terminate
  (Proposition~\ref{prop:shisout}).
\end{proof}

\section{Adequacy (Theorem~\ref{thm:adequacy})}\label{app:adequacy}
\begin{lemma}[Meaningless leaves]\label{lem:rootactive}
  $\outm(x) = \bot$ iff the behaviour at $x$ is the $\bot$ leaf.
\end{lemma}
\noindent
By Proposition~\ref{prop:outwd}, $\outm(x) = \bot$ iff the silent process from $x$ exposes no layer.
The graph's $\bot$ leaf is by Definition~\ref{def:reach} exactly such a state, so the identification is
immediate.

\begin{proof}[Proof of Theorem~\ref{thm:adequacy}]
  The structure map $\outm$ (Definition~\ref{def:out}) is a $\Beh$-coalgebra $X \to \Beh X$, $\Beh = MF$
  (Definition~\ref{def:behfun}); the graph $(R(M),\outm)$ is its restriction to the reachable states. By
  finality there is a unique $\Beh$-coalgebra morphism $\behav \colon X \to \nu\Beh$, determined layer
  for layer by $\outm$: $\behav(x)$ exposes the layer $\outm(x)$ with children $\behav$ of the successor
  states, or the $\bot$ leaf when $\outm(x) = \bot$ (Lemma~\ref{lem:rootactive}). Since
  $\outm = \iter{\stepm}$ runs the silent steps of $\stepm$ (Proposition~\ref{prop:outwd}), $\behav$ is
  the behaviour of $\stepm$ up to silent steps, and the graph $(R(M),\outm)$, whose unfolding is
  $\behav(M)$, is its finite presentation.
\end{proof}

\section{Termination (Theorem~\ref{thm:termination})}\label{app:termination}
\begin{proof}[Proof of Theorem~\ref{thm:termination}]
  Tabling descends only through reachable states (Definition~\ref{def:reach}): the children of an
  exposed layer and, via the silent steps of $\outm$, the silent successors. State equality is decidable,
  so the tabling key is testable.

  \emph{Finite $\Rightarrow$ halts.} If $R(M)$ is finite, tabling $\outm$ visits a finite set and halts:
  the silent collapse $\outm(x)$ saturates each finite silent orbit, and the reachable closure saturates
  $R(M)$.

  \emph{Halts $\Rightarrow$ finite.} Contrapositive: if $R(M)$ is infinite, the descent visits
  infinitely many distinct states and the tabling never saturates.

  \emph{Graph rational.} When $R(M)$ is finite, $(R(M),\outm)$ is a finite graph; its unfolding has
  subtrees $\{\behav(y) : y \in R(M)\}$, finite, so it lies in $\rfp\Beh$
  (Definition~\ref{def:rational}).

  \emph{Two conditions.} The two halts above are independent (Remark~\ref{rem:tworationalities}): the
  silent orbit of a state may be infinite while $R(M)$ is finite, or conversely.
\end{proof}

% =====================================================================
\section{The Rational Fragment as a Domain-Specific Language}\label{app:dsl}
The terminating fragment of the construction, the rational behaviours, is exactly the world of
finite-state, regular, and least-fixpoint-over-a-finite-lattice computation, and a range of
practical domains live there. Each domain below is a pure $\lambda$-term whose behaviour is
rational, so the interpreter halts, and each is reproduced as a test\anon[~(supplementary
material)]{~\cite{colambda}}. Two reusable shapes recur: a least fixpoint over a Church
tuple of booleans (the Datalog engine of Section~\ref{sec:application}), and a recursion whose
subproblems are the nodes of a first-order datum, memoised for free because identical subproblems
are one interned state. Each claim is confined to a domain's rational core: regular lexing but not
general context-free parsing, rational-tree unification but not a general most-general unifier,
finite-state model checking but not infinite-state. The terms below are the actual encodings; the
Booleans are Church, $\mathtt{T} = \lambda a.\lambda b.\,a$ and $\mathtt{F} = \lambda a.\lambda b.\,b$,
with $\mathtt{or} = \lambda p.\lambda q.\,p\,\mathtt{T}\,q$ and
$\mathtt{and} = \lambda p.\lambda q.\,p\,q\,\mathtt{F}$, and $\mathtt{cons}$ and $Y$ are as in
Section~\ref{sec:bridge}.

\paragraph{Dynamic programming, memoisation for free.}
The decisive capability is memoisation (Section~\ref{sec:memo}): a recursion whose state space is a
tree or graph computes each distinct subproblem once, because a structurally repeated subtree, or a
graph node revisited along several paths, is one interned state. A perfect binary tree of depth $n$
whose nodes share both children is a graph of $n+1$ states that unfolds to $2^{n}$ leaves; a tree
DP over it is read in time linear in $n$, where a naive recursion would visit $2^{n}$ leaves. No
memo table is written; tabling on first-order state identity is the memoisation, which the pure
$\lambda$-calculus lacks without a decidable identity on subproblems, value equality of closures
being undecidable.
\[
  \mathtt{node} = \lambda l.\lambda r.\lambda m.\lambda f.\,m\,l\,r, \qquad
  \mathtt{leaf} = \lambda v.\lambda m.\lambda f.\,f\,v,
\]
\[
  \mathtt{any} = Y\,\big(\lambda\mathit{self}.\lambda t.\;
  t\,(\lambda l.\lambda r.\,\mathtt{or}\,(\mathit{self}\,l)\,(\mathit{self}\,r))\,(\lambda v.\,v)\big),
  \qquad t_0 = \mathtt{leaf}\;\mathtt{F},\quad t_{k+1} = \mathtt{node}\;t_k\;t_k,
\]
so $\mathtt{any}\,t_n$ reads to $\mathtt{F}$ with each of the $n+1$ shared subtrees evaluated once.

\paragraph{Game search with a transposition table.}
Minimax is an and-or tree: a maximising position takes the disjunction over its moves, a minimising
position the conjunction, and a terminal position a Boolean. A \emph{transposition}, the same
position reached by a different move order, is the same interned node, so its value is computed
once. Interning is the transposition table of a game-playing program, with no table written by hand.
\[
  \mathtt{max} = \lambda l.\lambda r.\lambda x.\lambda n.\lambda f.\,x\,l\,r, \quad
  \mathtt{min} = \lambda l.\lambda r.\lambda x.\lambda n.\lambda f.\,n\,l\,r, \quad
  \mathtt{leaf} = \lambda v.\lambda x.\lambda n.\lambda f.\,f\,v,
\]
\[
  \mathtt{value} = Y\,\big(\lambda\mathit{self}.\lambda p.\;
    p\,(\lambda l.\lambda r.\,\mathtt{or}\,(\mathit{self}\,l)(\mathit{self}\,r))\,
  (\lambda l.\lambda r.\,\mathtt{and}\,(\mathit{self}\,l)(\mathit{self}\,r))\,(\lambda v.\,v)\big).
\]

\paragraph{Graph reachability and program analysis.}
Reachability and transitive closure over a directed graph, including one with cycles, are the least
fixpoint of the immediate-consequence operator, terminating because the lattice is finite; the same
shape decides reachability of a state in a finite transition system. Andersen-style points-to (alias)
analysis \citeTODO{andersen, doop}, in which $\mathtt{pointsTo}(p,o)$ follows from
$\mathtt{assign}(p,q)$ and $\mathtt{pointsTo}(q,o)$, is monotone Datalog, and a compiler reads it as
the same boolean fixpoint. The cyclic graph $a\to b\to c\to a$ with $c\to d$, and the points-to
program for $a=\mathtt{new}\,o_1;\,b=a;\,c=b$, are the clause sets
\[
  \begin{aligned}
    &\mathtt{reach}(a),\quad \mathtt{reach}(b)\leftarrow\mathtt{reach}(a),\quad
    \mathtt{reach}(c)\leftarrow\mathtt{reach}(b),\\
    &\mathtt{reach}(a)\leftarrow\mathtt{reach}(c),\quad \mathtt{reach}(d)\leftarrow\mathtt{reach}(c);\\
    &\mathtt{pt}(a,o_1),\quad \mathtt{pt}(b,X)\leftarrow\mathtt{pt}(a,X),\quad
    \mathtt{pt}(c,X)\leftarrow\mathtt{pt}(b,X),
  \end{aligned}
\]
so $\mathtt{reach}(d)$ and $\mathtt{pt}(c,o_1)$ hold while $\mathtt{reach}(e)$ and $\mathtt{pt}(c,o_2)$
do not.

\paragraph{Lexical analysis.}
A deterministic finite automaton is a finite cyclic transition graph, and running it is a fold over
the input that threads the state. Regular languages are the archetypal rational behaviour, so a lexer
is a $\lambda$-term whose behaviour is rational; the name \emph{rational} is itself the regular, or
Kleene, sense (Definition~\ref{def:rational}). Encoding the even state and the symbol $a$ as
$\mathtt{T}$, and the odd state and $b$ as $\mathtt{F}$, both two-way selectors,
\[
  \delta = \lambda s.\lambda x.\,s\,(x\,\mathtt{F}\,\mathtt{T})\,(x\,\mathtt{T}\,\mathtt{F}), \qquad
  \mathtt{accepts} = \lambda s.\,s\,\mathtt{T}\,\mathtt{F},
\]
\[
  \mathtt{run} = Y\,\big(\lambda\mathit{self}.\lambda s.\lambda w.\;
  w\,(\lambda h.\lambda t.\,\mathit{self}\,(\delta\,s\,h)\,t)\,s\big),
\]
and $\mathtt{accepts}\,(\mathtt{run}\,\mathtt{T}\,w)$ is $\mathtt{T}$ exactly when $w$ has an even
number of $a$s.

\paragraph{Unification and recursive types.}
Equality of two rational (cyclic) terms is decided by their behaviour through cycle-detecting
bisimulation (Corollary~\ref{cor:decidable}). This is the decidable core shared by rational-tree
unification \citeTODO{colmerauer} and equirecursive ($\mu$-) type equality
\citeTODO{amadio-cardelli, brandt-henglein}: two constructions of one rational term read out
identically, and a different term reads out differently. For instance $Y\,(\mathtt{cons}\;0)$ and
$Y\,(\lambda s.\,\mathtt{cons}\;0\;s)$ read out identically, while $Y\,(\mathtt{cons}\;1)$ differs.

\paragraph{The metalanguage needs tabling too.}
The interpreter's own host-language term builder illustrates the point reflexively. Writing a shared
subterm by reusing a builder gives a term whose construction unfolds the sharing exponentially, and
a cyclic builder would not terminate, until the builder is memoised by binder depth, the
host-language analogue of tabling. Intuitive code is slow or non-terminating exactly where the
metalanguage does not table. The memo key here is the syntactic builder rather than the first-order
behavioural identity the calculus folds on, but the moral is the same.

\paragraph{Further instances.}
The same rational core underlies others not reproduced here: equality saturation over e-graphs
\citeTODO{egg, egglog} and binary decision diagrams \citeTODO{bryant}, both of which are interning
of a shared graph; garbage-collection reachability over a cyclic heap; abstract interpretation as a
fixpoint over a finite abstract domain; and, with a non-trivial effect monad, the stationary
distribution of a finite Markov chain. Each is a finite-state, regular, or finite-lattice-fixpoint
computation, hence a rational behaviour the construction reads.

% =====================================================================
\section{Compilation to Python by Analysis-Driven Specialization}\label{app:compiler}
The interpreter is the default evaluator, and it is already a fixpoint-thunk graph: every node's
weak head normal form is a \texttt{fixpoint\_cached\_property} thunk, interned and possibly cyclic,
so the term graph is the thunk graph the interpreter folds, and handing that graph back is the
identity. A compiler can do something different only by choosing a different evaluation strategy, and
it preserves the interpreter's result only when an analysis certifies the choice is sound. The
accompanying implementation\anon[~(supplementary material)]{~\cite{colambda}} carries this
out as a self-hosting compiler from the pure $\lambda$-calculus to a Python AST that specializes a
sub-term to a strict or a call-by-name Python target when an analysis certifies the result is
unchanged, and leaves everything else to the interpreter. Specialization is local, the analysis is
itself a $\lambda$-term, and a depth bound turns reduction into a family of readings between the raw
term and the weak head normal form. The pieces below are executable, each reproduced as a test.

\paragraph{The target: a Scott-encoded Python AST.}
The target is a Python AST represented in the pure $\lambda$-calculus by Scott encoding, a value of
an $n$-constructor type being $\lambda h_0 \dots h_{n-1}.\,h_{\mathit{tag}}\,\mathit{fields}$. The
encoding and its decoder are generated by reflection on Python's \texttt{ast} module: the supported
node classes give the constructors, each class's fields give the arity and order, and a field
carries a kind tag (a child, a list, an integer, a string as character codes). Real Python
round-trips through this encoding, a function with a \texttt{while} loop included, and the decoded
program runs.

\paragraph{The compiler, in the $\lambda$-calculus.}
A quoted source term is a Scott value over $\mathtt{QVar}\,i$, $\mathtt{QLam}\,b$,
$\mathtt{QApp}\,f\,a$ (de Bruijn). The compiler is the pure term
\[
  \mathtt{compile} = Y\,\big(\lambda\mathit{self}.\lambda d.\lambda q.\;
    q\,(\lambda i.\,\mathtt{PyVar}\,(d - 1 - i))\,
    (\lambda b.\,\mathtt{PyLam}\,d\,(\mathit{self}\,(d{+}1)\,b))\,
  (\lambda f.\lambda a.\,\mathtt{PyApp}\,(\mathit{self}\,d\,f)\,(\mathit{self}\,d\,a))\big),
\]
mapping an abstraction to a Python \texttt{lambda}, an application to a call, and a variable to a
name; it threads the binder depth $d$ so a de Bruijn index $i$ becomes the level $d - 1 - i$
(Church subtraction), giving stable parameter names. Compiled identity, $K$, Church numerals, and
arithmetic applications execute in Python and compute what the source terms compute, as the generated
comparison under \emph{Compiled examples} below shows. The compiler
emits one of three targets, selected by the analysis below: a strict (eager) target; a call-by-name
lazy target; and a fixpoint target that is the interpreter itself, since the fixpoint-thunk graph is
the AST. The eager and lazy targets differ from the interpreter only in evaluation strategy, so they
are sound only under the conditions the analysis certifies; the lazy target uses a thunk forced on
reference, the call-by-name semantics matching weak-head reduction, so every normalizing term
computes its value, factorial and Fibonacci among them. The prelude arithmetic is cross-checked
against the interpreter on the lazy target.

\paragraph{Compiled examples.}
Each example below shows the specializer's output: a closed simply-typable term certified to the
strict call-by-value target, an untypable but normalizing self-application to call-by-need, and a
fold-requiring cyclic term left to the interpreter with its closed element compiled to a call-by-value
island. The flagship is the compiler itself. Untypable as a whole, since its fixpoint combinator
self-applies, it stays interpreted, but the specializer finds its maximal closed simply-typable
sub-terms, the strongly-normalizing combinators, and compiles each to a strict island, carving
call-by-value regions out of an otherwise interpreted program. A bound variable is named by a list
of naturals the compiler emits (the binder depth for the expression targets, the binder's
syntax-tree path for call-by-need) and a single decoder, shared by every target, renders it as the
Python identifier \texttt{v} with an underscore-separated integer per segment, so distinct binders
get distinct names by construction. The listings are generated from the compiler, not written by
hand, and the compiler's staged self-compilations are committed alongside as generated modules.
% OUTDATED: the surrounding prose describes the old specializer (islands / three targets).
% generated/compiler-examples.tex is now produced by the DEFUN compiler instead: each example
% term is a module of @interned closure classes (regenerate with co-lambda-defun-examples).
% Generated by co_lambda_examples._examples (co-lambda-defun-examples). Do not edit;
% regenerate with: python -m co_lambda_examples._examples

\medskip\noindent\textbf{Identity.}\quad $\lambda x.\, x$
\begin{lstlisting}
@interned
class vg_90808ac1cd37d6ee:

    def __call__(self, a):
        return a
compiled = vg_90808ac1cd37d6ee()
\end{lstlisting}

\medskip\noindent\textbf{The constant combinator $K$.}\quad $\lambda x.\, \lambda y.\, x$
\begin{lstlisting}
@interned
class vg_16f2ce6352e69203:

    def __call__(self, a):
        return vg_a3d71a2af140cd9a(a)

@interned
class vg_a3d71a2af140cd9a:
    cap_0: Lambda

    def __call__(self, a):
        return self.cap_0
compiled = vg_16f2ce6352e69203()
\end{lstlisting}

\medskip\noindent\textbf{Church numeral $2$.}\quad $\lambda x.\, \lambda y.\, x\, (x\, y)$
\begin{lstlisting}
@interned
class vg_e15d7ee56e02d808:

    def __call__(self, a):
        return vg_e874ce9954e7e0b9(a)

@interned
class vg_e874ce9954e7e0b9:
    cap_0: Lambda

    def __call__(self, a):
        return Thunk(self.cap_0, Thunk(self.cap_0, a))
compiled = vg_e15d7ee56e02d808()
\end{lstlisting}

\medskip\noindent\textbf{Successor.}\quad $\lambda x.\, \lambda y.\, \lambda z.\, y\, (x\, y\, z)$
\begin{lstlisting}
@interned
class vg_2f5ca52fef3a44f2:

    def __call__(self, a):
        return vg_9a8bc0dc64816bfe(a)

@interned
class vg_9a8bc0dc64816bfe:
    cap_0: Lambda

    def __call__(self, a):
        return vg_a3db8e47dd43cac9(a, self.cap_0)

@interned
class vg_a3db8e47dd43cac9:
    cap_0: Lambda
    cap_1: Lambda

    def __call__(self, a):
        return Thunk(self.cap_0, Thunk(Thunk(self.cap_1, self.cap_0), a))
compiled = vg_2f5ca52fef3a44f2()
\end{lstlisting}

\medskip\noindent\textbf{Addition.}\quad $\lambda x.\, \lambda y.\, \lambda z.\, \lambda w.\, x\, z\, (y\, z\, w)$
\begin{lstlisting}
@interned
class vg_99cfd8ca8f1d742b:
    cap_0: Lambda
    cap_1: Lambda

    def __call__(self, a):
        return vg_cdf25b0daf7dd82a(self.cap_0, a, self.cap_1)

@interned
class vg_bb5f1aa674606fcf:
    cap_0: Lambda

    def __call__(self, a):
        return vg_99cfd8ca8f1d742b(self.cap_0, a)

@interned
class vg_cdf25b0daf7dd82a:
    cap_0: Lambda
    cap_1: Lambda
    cap_2: Lambda

    def __call__(self, a):
        return Thunk(Thunk(self.cap_0, self.cap_1), Thunk(Thunk(self.cap_2, self.cap_1), a))

@interned
class vg_d96a8b24aabe3dee:

    def __call__(self, a):
        return vg_bb5f1aa674606fcf(a)
compiled = vg_d96a8b24aabe3dee()
\end{lstlisting}

\medskip\noindent\textbf{Multiplication.}\quad $\lambda x.\, \lambda y.\, \lambda z.\, x\, (y\, z)$
\begin{lstlisting}
@interned
class vg_1f95eda5152a1edd:

    def __call__(self, a):
        return vg_ac4715584220052b(a)

@interned
class vg_ac4715584220052b:
    cap_0: Lambda

    def __call__(self, a):
        return vg_f2c03833bbc18fa7(self.cap_0, a)

@interned
class vg_f2c03833bbc18fa7:
    cap_0: Lambda
    cap_1: Lambda

    def __call__(self, a):
        return Thunk(self.cap_0, Thunk(self.cap_1, a))
compiled = vg_1f95eda5152a1edd()
\end{lstlisting}

\medskip\noindent\textbf{Zero test.}\quad $\lambda x.\, x\, (\lambda y.\, \lambda z.\, \lambda w.\, w)\, (\lambda y.\, \lambda z.\, y)$
\begin{lstlisting}
@interned
class vg_16f2ce6352e69203:

    def __call__(self, a):
        return vg_a3d71a2af140cd9a(a)

@interned
class vg_573dd662ffb3725e:

    def __call__(self, a):
        return vg_5e82a4badd71c4a6()

@interned
class vg_5e82a4badd71c4a6:

    def __call__(self, a):
        return vg_90808ac1cd37d6ee()

@interned
class vg_90808ac1cd37d6ee:

    def __call__(self, a):
        return a

@interned
class vg_a3d71a2af140cd9a:
    cap_0: Lambda

    def __call__(self, a):
        return self.cap_0

@interned
class vg_bfb29603bf3e65f5:

    def __call__(self, a):
        return Thunk(Thunk(a, vg_573dd662ffb3725e()), vg_16f2ce6352e69203())
compiled = vg_bfb29603bf3e65f5()
\end{lstlisting}

\medskip\noindent\textbf{Self-application.}\quad $\lambda x.\, x\, x$
\begin{lstlisting}
@interned
class vg_99b73b6fd3d9f13d:

    def __call__(self, a):
        return Thunk(a, a)
compiled = vg_99b73b6fd3d9f13d()
\end{lstlisting}


\paragraph{The soundness analysis.}
A static analysis chooses the target. A term that is simply typable, decided by algorithm-W
inference with an occurs check so the self-application $x\,x$ of $Y$ and $\Omega$ fails, is
strongly normalizing, so the strict eager target terminates with the interpreter's normal form.
Otherwise the interpreter serves as a sound oracle: it reads the behaviour out, and if the behaviour
is a finite normal form (no folded back edge and no $\bot$), the term is normalizing and the
call-by-name lazy target reaches the same value. Failing both, the term keeps the fixpoint default,
the interpreter, which folds every rational behaviour correctly. Church arithmetic compiles eager;
factorial and Fibonacci applied to a numeral compile lazy; $Y\,(\mathtt{cons}\;0)$ and $\Omega$ stay
interpreted. Simple typability is a sound but conservative certificate, and normalization is
undecidable, so the lazy test runs the safe interpreter and reads off whether it folded; no totality
is claimed, and anything the analysis does not certify stays interpreted.

\paragraph{The analysis, in the $\lambda$-calculus.}
The certificate that selects what to specialize is itself a pure $\lambda$-term, run by the
interpreter on the quoted program, so the calculus analyzes its own programs. The first and simplest
is a closedness check
\[
  \mathtt{closed} = Y\,\big(\lambda\mathit{self}.\lambda d.\lambda q.\;
    q\,(\lambda i.\, i < d)\,
    (\lambda b.\,\mathit{self}\,(d{+}1)\,b)\,
  (\lambda f.\lambda a.\,\mathit{and}\,(\mathit{self}\,d\,f)\,(\mathit{self}\,d\,a))\big),
\]
recursing structurally with $Y$ and threading a Church-numeral binder depth: a variable is in scope
when its index is below the depth, an abstraction recurses at $d+1$, an application is closed when
both sides are, with $<$ and $\mathit{and}$ the usual Church encodings. Closedness is exactly the
precondition a compiled island needs, and richer certificates (rationality by folding, typability)
layer on in the same style. That the analysis is a $\lambda$-term rather than host code is the point:
tabling-based reduction expresses program analysis, not only evaluation.

\paragraph{Local specialization through a foreign-function node.}
Specialization is local, not whole-program: a single node kind embeds a compiled Python callable in
the term graph, taking argument nodes to a result node, so the node graph is the interoperation
boundary. A closed sub-term the analysis certifies is compiled once to such a node and spliced into an
otherwise interpreted graph; the interpreter drives the island and folds around it. In
$Y\,(\mathtt{cons}\,(3 \times 3))$ the cyclic $\mathtt{cons}$ stays interpreted while the closed
element compiles to an island, and the readout, a cyclic stream of $9$, is unchanged. A solution
written once is thus compiled once and run on many $\lambda$-term inputs from the host, inputs and
results staying $\lambda$ values.

\paragraph{Bootstrap.}
The fixpoint is the ordinary combinator $Y$: its self-application is untypable and has no finite
normal form, so the analyses never send the recursion to a compiled target and it stays
interpreted. The compiler then compiles
its own source to a Python program, and that self-compiled compiler, evaluated as a Python function
and run on host-encoded source terms, produces the same output as the original and itself runs.

\paragraph{A weak head and combinatory fusion: depth-bounded reduction.}
The compiler's reference semantics fuses two reductions. A combinatory (SK) reduction pushes a
substitution one constructor inward, as an explicit substitution calculus
does~\citeTODO{abadi-curien-levy, benaissa-lescanne}, instead of copying the whole spine, which is
efficient, but its $S$ rule, $(\lambda a.\,E_1\,E_2)\,b \to ((\lambda a.\,E_1)\,b)\,((\lambda a.\,E_2)\,b)$,
turns the argument into suspensions that grow on every turn of a cycle (the sharing problem of
combinator graph reduction~\citeTODO{wadsworth-graph, lamping}), so the cycle never re-forms one
interned node and the fold never fires; reducing the suspensions first (call by value) re-forms the node but
diverges on a productive cycle such as $Y\,(\mathtt{cons}\;0)$. The construction instead takes
call-by-name beta, where the argument is inserted by reference rather than reduced, so tabling folds
both $\Omega$ (to $\bot$) and $Y\,(\mathtt{cons}\;0)$ (to a finite cyclic graph), together with a depth
bound on the substitution that leaves a still-unfired redex $(\lambda a.\,T)\,a$ as a guarded let-stub
rather than substituting the whole, possibly infinite, tree. Because substitution returns a closed
subterm by reference, closed cyclic data is shared and only the open redex-body spine is copied, the
minimal traversal the fold requires. Bounding the reduction gives a family of structure maps between
the raw term ($d = 0$) and the weak head normal form ($d$ large), with $d = 1$ the one-layer reading;
growing $d$ climbs the $\bot$-completed approximation order~\citeTODO{barendregt, levy-longo}, whose
limit is the Levy-Longo reading. So
the semantics folds the common cycles, as weak head normalization does, while never doing the
unbounded whole-tree substitution that motivates combinatory reduction.

\paragraph{The halting boundary.}
Faithfulness is convergence to the unique fixpoint, not structural identity: since the least and
greatest fixpoints coincide (Theorem~\ref{thm:unique}), an operational strategy is faithful as long
as it reaches that fixpoint rather than diverging, so a folded cycle may be unrolled into a finite
tree or reshaped into a different cycle and stay faithful, and the only unfaithful move is unrolling a
cycle into an infinite tree. This places compilation on a halting spectrum. Plain lazy reduction folds
nothing and halts only on a finite normal form. The interpreter folds at every depth and halts on
exactly the rational behaviour, the optimal boundary (Theorems~\ref{thm:termination},
\ref{thm:optimal}). Bounded compilation lies strictly between: every contraction it fires is interned
and its re-entry folded, so a cycle that closes within the depth folds and halts, and it halts on more
than plain lazy reduction; a cycle that closes only deeper is left suspended, and forcing it past the
bound without returning to the interpreter is where compiled code can diverge.

\paragraph{Loops from folding.}
A code generator
$\mathtt{gen} = Y\,\big(\lambda\mathit{self}.\lambda s.\;
s\,(\lambda h.\lambda t.\,\mathtt{Yield}\,h\,(\mathit{self}\,t))\,\mathtt{Stop}\big)$
maps a stream to a quoted output stream, with nothing loop-aware in it. Run on a cyclic source such
as $Y\,(\mathtt{cons}\;0)$, the recursive $\mathit{self}\,t$ re-enters the same interned state, so
the interpreter folds the back edge and the output is a cyclic quoted stream; a cycle-aware decoder
turns it into a Python generator with a \texttt{while} loop, which runs and yields the stream. On a
finite source the same $\mathtt{gen}$ yields a finite, loopless generator. The loop in the compiled
program is exactly the interpreter's fold of the rational behaviour made imperative, with no sharing
test in the $\lambda$-term. This is a further specialization, beyond the eager and lazy targets, of a
productive stream whose behaviour is rational.

\paragraph{Scope.}
The pieces are executable and tested. Read as Futamura projections \citeTODO{futamura},
specializing a program under a soundness analysis is an ahead-of-time compiler, and the fold that
decides $\Omega$ as $\bot$ and folds a cyclic datum is the interpreter's default the analysis falls
back to. Meta-tracing does not apply here: a meta-tracing just-in-time compiler
\citeTODO{bolz-tracing} compiles traces of an interpreter, which requires the interpreter to be
written in the traced language, a $\lambda$ self-interpreter, which this construction does not use;
here the $\lambda$ term is the program and the interpreter folds its behaviour directly.

\section{Self-Compilation: The Multi-Stage Climb as a Matrix}\label{app:benchmark}
The compiler compiles itself, and the island depth is its option (Appendix~\ref{app:compiler}): a larger
depth compiles more of the compiler as native islands and leaves less to the interpreter. Each closed
sub-term is certified one of three ways: a simply-typable sub-term becomes a strict call-by-value island,
a normalizing but untypable sub-term becomes a lazy island (call-by-need or call-by-name), and the rest
stays interpreted. A progressive bootstrap tests whether that makes the compiler faster: the compiler run
on the interpreter compiles the compiler into a small-island compiler; that compiler compiles the
compiler into a larger-island compiler; and so on, each stage compiled by the previous.
Table~\ref{tab:self-compilation} expands the climb into a matrix (every compiler version compiling the
compiler's own source at every target depth, each in a fresh process, reporting wall-clock time and peak
resident memory), doubled across the two lazy-island regimes (call-by-need memoises shared sub-results,
call-by-name recomputes them). The matrix is shown to target depth 64; the depth 128 and 192 columns each
need several gigabytes per cell, so they are not run here, though the staged compilers are committed for
every depth (the depth-192 compiler carries the two deep combinators below as whole islands).

\begin{table}[h]
\centering
% OUTDATED: this multi-stage island-depth matrix belongs to the removed specializer. The current
% self-compilation measurement is generated/defun-benchmark.tex (interpreter vs defunctionalized,
% with the DEFUN-on-DEFUN bootstrap row); regenerate with co-lambda-defun-benchmark.
% Generated by co_lambda._benchmark (co-lambda-benchmark). Do not edit.
% Matrix doubled across lazy regimes; rows = engine version, columns = target island size.
% Each regime block: wall-clock time (s) above, peak resident memory (GB) below.
\medskip\noindent\textbf{Call-by-need lazy islands (memoise).}\par\smallskip
\begin{tabular}{lrrrr}
\hline
Engine & island 0 & island 8 & island 32 & island 64 \\
\hline
  interpreter & 20.8 & 45.1 & 67.9 & 77.4 \\
  island 0 & 20.7 & 48.1 & 66.3 & 77.9 \\
  island 8 & 20.8 & 46.4 & 67.2 & 76.7 \\
  island 32 & 20.3 & 45.9 & 65.4 & 77.8 \\
  island 64 & 20.5 & 45.7 & 65.5 & 77.6 \\
\hline
\end{tabular}
\par\smallskip
\begin{tabular}{lrrrr}
\hline
Engine & island 0 & island 8 & island 32 & island 64 \\
\hline
  interpreter & 0.71 & 1.86 & 2.85 & 3.29 \\
  island 0 & 0.71 & 1.86 & 2.85 & 3.29 \\
  island 8 & 0.71 & 1.86 & 2.86 & 3.29 \\
  island 32 & 0.71 & 1.86 & 2.86 & 3.30 \\
  island 64 & 0.71 & 1.86 & 2.86 & 3.30 \\
\hline
\end{tabular}
\par
\medskip\noindent\textbf{Call-by-name lazy islands (recompute).}\par\smallskip
\begin{tabular}{lrrrr}
\hline
Engine & island 0 & island 8 & island 32 & island 64 \\
\hline
  interpreter & 20.8 & 45.1 & 67.9 & 77.4 \\
  island 0 & 20.7 & 48.1 & 66.3 & 77.9 \\
  island 8 & 20.8 & 46.4 & 67.2 & 76.7 \\
  island 32 & 20.4 & 45.4 & 65.2 & 76.9 \\
  island 64 & 20.6 & 45.4 & 65.4 & 78.4 \\
\hline
\end{tabular}
\par\smallskip
\begin{tabular}{lrrrr}
\hline
Engine & island 0 & island 8 & island 32 & island 64 \\
\hline
  interpreter & 0.71 & 1.86 & 2.85 & 3.29 \\
  island 0 & 0.71 & 1.86 & 2.85 & 3.29 \\
  island 8 & 0.71 & 1.86 & 2.86 & 3.29 \\
  island 32 & 0.71 & 1.86 & 2.86 & 3.30 \\
  island 64 & 0.71 & 1.86 & 2.86 & 3.30 \\
\hline
\end{tabular}
\par

\caption{The multi-stage climb as a matrix, doubled across the two lazy-island regimes (call-by-need
memoise and call-by-name recompute): rows are the compiler version doing the compiling (the interpreter,
then the compiler self-compiled at each island depth), columns the target island depth. Within each
regime block the upper tabular is wall-clock time in seconds and the lower is peak resident memory in
gigabytes, each of one fresh process compiling the compiler's own source. The linear bootstrap is the
diagonal.}
\label{tab:self-compilation}
\end{table}

The matrix is flat down every column, and the two regimes coincide: for a fixed target depth the
interpreter and every self-compiled engine take essentially the same time and the same memory, and
call-by-need and call-by-name agree. So the cost is set by the target depth alone, not by which version
of the compiler does the compiling nor by the lazy regime. The reason is the soundness discipline the
islands obey: a compiled island must converge exactly where the interpreter does. The compiler's hot path
is its analysis, the typability and normalization checks that decide each sub-term's island, driven by
the same $Y$ fixpoint that makes the compiler untypable. As a value that recursion has no normal form, so
it can be neither a call-by-value island (which requires strong normalization) nor a lazy island (which
requires a finite normal form); it stays interpreted in every engine. Its per-step body is a different
matter: that body is a closed term that uses the recursive call only as a parameter (no self-application),
so several of the analysis step bodies are simply typable and are compiled as islands. They still do not
change the cost, and measuring a reflect boundary shows why. A typable step body compiled to native code
runs with no interpreter involvement when its arguments are host values. In the hot path, though, its
arguments arrive from the untypable driver as interpreter nodes, and the node-to-host bridge a positive
arity island would need so the step body can branch on them is itself interpreter-driven: it forces each
sub-result back through the interpreter (one application per step, so interpreter speed), and a variant
that does not force diverges under read-back. So the per-step work is performed by the interpreter either
way; the native code helps only where the driver and the data are already native, which the untypable
recursion driver is not. That, not the size of the islands, is why the matrix is flat and the two regimes
coincide. The cost grows only with the target depth, because a deeper target types deeper sub-terms.
The deeper obstacle is that soundness and performance are not the same kind of property here. That a
compiled island converges exactly where the interpreter does is a local fact, provable for one island in
isolation. Performance is global: the interpreter's tabling is a single cache shared across the whole
compilation, so carving a recursion into an island, which keeps its own memo, gives up the cross-term
sharing the analysis was built on (the closedness measure, for instance, is deliberately depth-free so
the interner shares it at every position and a whole-tree scan is linear). Removing one such computation
from the shared cache is sound but can slow code arbitrarily far away that depended on the same entry,
and no local analysis can rule that out. A sound speedup of the hot path therefore cannot be licensed by
a static local certificate; it would need a feedback-directed heuristic, in the spirit of hardware branch
prediction or a tracing just-in-time compiler, that uses an observed execution profile to predict which
regions repay native compilation and leaves the rest to the shared interpreter, correcting itself when a
prediction does not pay. Speeding up the shared interpreter itself is the complementary alternative whose
benefit is global by construction and so cannot regress a distant cache.

% =====================================================================
%% The body cites via \citeTODO (red draft markers), never \cite, so
%% \nocite{*} forces every references.bib entry into the printed list,
%% reproducing the former inline thebibliography under ACM formatting.
\nocite{*}
\bibliographystyle{ACM-Reference-Format}
\bibliography{references}

\end{document}
